% Options for packages loaded elsewhere
\PassOptionsToPackage{unicode}{hyperref}
\PassOptionsToPackage{hyphens}{url}
\PassOptionsToPackage{dvipsnames,svgnames,x11names}{xcolor}
\documentclass[
  english,
  11pt,
  reqno]{amsart}
\usepackage{xcolor}
\usepackage{amsmath,amssymb}
\setcounter{secnumdepth}{-\maxdimen} % remove section numbering
\usepackage{iftex}
\ifPDFTeX
  \usepackage[T1]{fontenc}
  \usepackage[utf8]{inputenc}
  \usepackage{textcomp} % provide euro and other symbols
\else % if luatex or xetex
  \usepackage{unicode-math} % this also loads fontspec
  \defaultfontfeatures{Scale=MatchLowercase}
  \defaultfontfeatures[\rmfamily]{Ligatures=TeX,Scale=1}
\fi
\usepackage{lmodern}
\ifPDFTeX\else
  % xetex/luatex font selection
\fi
% Use upquote if available, for straight quotes in verbatim environments
\IfFileExists{upquote.sty}{\usepackage{upquote}}{}
\IfFileExists{microtype.sty}{% use microtype if available
  \usepackage[]{microtype}
  \UseMicrotypeSet[protrusion]{basicmath} % disable protrusion for tt fonts
}{}
\usepackage{longtable,booktabs,array}
\newcounter{none} % for unnumbered tables
\usepackage{calc} % for calculating minipage widths
% Correct order of tables after \paragraph or \subparagraph
\usepackage{etoolbox}
\makeatletter
\patchcmd\longtable{\par}{\if@noskipsec\mbox{}\fi\par}{}{}
\makeatother
% Allow footnotes in longtable head/foot
\IfFileExists{footnotehyper.sty}{\usepackage{footnotehyper}}{\usepackage{footnote}}
\makesavenoteenv{longtable}
% definitions for citeproc citations
\NewDocumentCommand\citeproctext{}{}
\NewDocumentCommand\citeproc{mm}{%
  \begingroup\def\citeproctext{#2}\cite{#1}\endgroup}
\makeatletter
 % allow citations to break across lines
 \let\@cite@ofmt\@firstofone
 % avoid brackets around text for \cite:
 \def\@biblabel#1{}
 \def\@cite#1#2{{#1\if@tempswa , #2\fi}}
\makeatother
\newlength{\cslhangindent}
\setlength{\cslhangindent}{1.5em}
\newlength{\csllabelwidth}
\setlength{\csllabelwidth}{3em}
\newenvironment{CSLReferences}[2] % #1 hanging-indent, #2 entry-spacing
 {\begin{list}{}{%
  \setlength{\itemindent}{0pt}
  \setlength{\leftmargin}{0pt}
  \setlength{\parsep}{0pt}
  % turn on hanging indent if param 1 is 1
  \ifodd #1
   \setlength{\leftmargin}{\cslhangindent}
   \setlength{\itemindent}{-1\cslhangindent}
  \fi
  % set entry spacing
  \setlength{\itemsep}{#2\baselineskip}}}
 {\end{list}}
\usepackage{calc}
\newcommand{\CSLBlock}[1]{\hfill\break\parbox[t]{\linewidth}{\strut\ignorespaces#1\strut}}
\newcommand{\CSLLeftMargin}[1]{\parbox[t]{\csllabelwidth}{\strut#1\strut}}
\newcommand{\CSLRightInline}[1]{\parbox[t]{\linewidth - \csllabelwidth}{\strut#1\strut}}
\newcommand{\CSLIndent}[1]{\hspace{\cslhangindent}#1}
\ifLuaTeX
\usepackage[bidi=basic,shorthands=off]{babel}
\else
\usepackage[bidi=default,shorthands=off]{babel}
\fi
\ifLuaTeX
  \usepackage{selnolig} % disable illegal ligatures
\fi
\setlength{\emergencystretch}{3em} % prevent overfull lines
\providecommand{\tightlist}{%
  \setlength{\itemsep}{0pt}\setlength{\parskip}{0pt}}
% Portable additions for the arXiv-oriented profile.
% Keep packages within the standard TeX Live distribution.
\usepackage{mathtools}
\usepackage{amssymb}
\usepackage{amsthm}
\usepackage{booktabs}
\usepackage{listings}
\usepackage{microtype}
\usepackage{centernot}
\usepackage{newunicodechar}

\ifPDFTeX
  % unicode-math supplies mathscr on Xe/LuaTeX; pdfTeX needs its own font.
  \usepackage{mathrsfs}
  % Explicit, meaning-preserving mappings only. Other unsupported input errors
  % are fatal; never discard characters or guess a different math expression.
  \newunicodechar{α}{\ensuremath{\alpha}}
  \newunicodechar{β}{\ensuremath{\beta}}
  \newunicodechar{γ}{\ensuremath{\gamma}}
  \newunicodechar{δ}{\ensuremath{\delta}}
  \newunicodechar{ε}{\ensuremath{\varepsilon}}
  \newunicodechar{θ}{\ensuremath{\theta}}
  \newunicodechar{λ}{\ensuremath{\lambda}}
  \newunicodechar{μ}{\ensuremath{\mu}}
  \newunicodechar{π}{\ensuremath{\pi}}
  \newunicodechar{σ}{\ensuremath{\sigma}}
  \newunicodechar{ω}{\ensuremath{\omega}}
  \newunicodechar{ℂ}{\ensuremath{\mathbb{C}}}
  \newunicodechar{ℕ}{\ensuremath{\mathbb{N}}}
  \newunicodechar{ℚ}{\ensuremath{\mathbb{Q}}}
  \newunicodechar{ℝ}{\ensuremath{\mathbb{R}}}
  \newunicodechar{ℤ}{\ensuremath{\mathbb{Z}}}
  \newunicodechar{∞}{\ensuremath{\infty}}
  \newunicodechar{∈}{\ensuremath{\in}}
  \newunicodechar{∅}{\ensuremath{\emptyset}}
  \newunicodechar{⊂}{\ensuremath{\subset}}
  \newunicodechar{⊆}{\ensuremath{\subseteq}}
  \newunicodechar{∩}{\ensuremath{\cap}}
  \newunicodechar{∪}{\ensuremath{\cup}}
  \newunicodechar{≤}{\ensuremath{\leq}}
  \newunicodechar{≥}{\ensuremath{\geq}}
  \newunicodechar{≠}{\ensuremath{\neq}}
  \newunicodechar{→}{\ensuremath{\rightarrow}}
  \newunicodechar{↦}{\ensuremath{\mapsto}}
  \newunicodechar{⇒}{\ensuremath{\Rightarrow}}
  \newunicodechar{⇔}{\ensuremath{\Leftrightarrow}}
  \newunicodechar{−}{\ensuremath{-}}
\fi
\newunicodechar{□}{\ensuremath{\square}}

% Portable negated hooked arrow used by several manuscripts.
\providecommand{\nhookrightarrow}{\centernot\hookrightarrow}

% U+220E END OF PROOF: render a small hollow square at the right margin.
\newcommand{\rightqedsquare}{%
  \ifhmode
    \unskip\nobreak\hfill\ensuremath{\square}%
  \else
    \noindent\hfill\ensuremath{\square}%
  \fi
}
\newunicodechar{∎}{\rightqedsquare}

% Reduce avoidable overfull lines without changing the source manuscript.
\setlength{\emergencystretch}{2em}
% Match amsart's theorem preskip while keeping Markdown statement bodies roman.
\makeatletter
\newcommand{\pandocstatementheading}[1]{%
  \par\addpenalty{\@beginparpenalty}%
  \addvspace{.5\baselineskip plus .2\baselineskip minus .2\baselineskip}%
  {\normalfont\bfseries #1\par}%
  \nobreak\@afterindenttrue\@afterheading
}
\makeatother
% House style: centered bold upright Abstract heading at abstract-body size.
% The native abstract box, margins and body font remain owned by amsart.
\usepackage{etoolbox}
\newcommand{\pandocabstractheading}{%
  \par\begingroup
    % Change only the heading style; keep the contents/bookmark title plain.
    \let\scshape\bfseries
    \section*{\abstractname}%
  \endgroup
  \nobreak\noindent\ignorespaces
}
\makeatletter
\patchcmd{\abstract}
  {\item[\hskip\labelsep\scshape\abstractname.]}
  {\item[]\pandocabstractheading}
  {}%
  {\PackageError{studio-ams-abstract}
     {Unsupported amsart abstract definition}
     {Check the installed class before updating the abstract-heading patch.}}
\makeatother
% Common LaTeX macros for CCG papers
% Keep this minimal and portable across engines.

% Math shortcuts
\newcommand{\N}{\mathbb{N}}
\newcommand{\Z}{\mathbb{Z}}
\newcommand{\Q}{\mathbb{Q}}
\newcommand{\R}{\mathbb{R}}
\newcommand{\C}{\mathbb{C}}
\newcommand{\eps}{\varepsilon}

% Theorem-like environments (lightweight)
% If you need full theoremstyles, consider amsthm or thmtools here.
% \usepackage{amsthm}   % already available via TeX Live if needed
% \newtheorem{theorem}{Theorem}
% \newtheorem{lemma}{Lemma}

% Code listings default (avoid minted to keep shell-escape off)
\lstset{basicstyle=\ttfamily\small,breaklines=true,columns=fullflexible}

\usepackage{bookmark}
\IfFileExists{xurl.sty}{\usepackage{xurl}}{} % add URL line breaks if available
\urlstyle{same}
\hypersetup{
  pdftitle={Local analytic Einstein metrics from free curvature data in four dimensions},
  pdfauthor={Lior Isthmus},
  pdflang={en},
  colorlinks=true,
  linkcolor={blue},
  filecolor={Maroon},
  citecolor={blue},
  urlcolor={blue},
  pdfcreator={LaTeX via pandoc}}

\title{Local analytic Einstein metrics from free curvature data in four
dimensions}
\author{Lior Isthmus\textsuperscript{1}}
\date{2026-11-18}

\usepackage{needspace}
\begin{document}
\raggedbottom
\begin{abstract}
We construct marked holomorphic Einstein four-germs from the two
complete chiral curvature records at a point. For every prescribed
numerical cosmological constant, a factorial growth bound on the free
records is necessary and sufficient for realization by such a germ. We
complete the ordered curvature jets, estimate the nonlinear horizontal
derivations, and integrate them on a scale of analytic record spaces.
The resulting radial curvature determines the metric through a Jacobi
equation. The metric, its inverse and its normal-coordinate realization
depend holomorphically on bounded record balls on common neighborhoods.
The same record flows represent short point and full spin-frame changes,
and the records fixed by Lorentz conjugation give the real-analytic
Lorentzian Einstein germs.
\end{abstract}
\maketitle


\markboth{\MakeUppercase{Lior Isthmus}}{\MakeUppercase{Local analytic Einstein metrics from free curvature\ldots{}}}

\protect\phantomsection\label{sec:es-abstract}

\tableofcontents

\footnotetext[1]{Email: \href{mailto:lioristhmus@proton.me}{\nolinkurl{lioristhmus@proton.me}}}
\setcounter{footnote}{1}

\protect\phantomsection\label{sec:es-introduction}

\subsection{1. Introduction}\label{introduction}

In four dimensions, a factorial growth bound on the free chiral
curvature derivatives at a point is necessary and sufficient for
realization by a marked holomorphic Einstein metric germ. The germ is
unique up to marked isometry, for every prescribed numerical
cosmological constant, including zero. The free data are the separately
symmetric covariant derivatives of the two Weyl spinors. Their
Lorentz-real form gives exactly the marked real-analytic Lorentzian
Einstein germs. The construction applies to every algebraic type of the
Weyl curvature.

On each bounded ball of analytic records, the metrics and their inverses
are bounded on a common normal-coordinate neighborhood and depend
holomorphically on the records. Bounded holomorphic families also have
common domains for the coordinate maps and radial spin frames. The
record flows describe changes of base point and full spin frame, with
common representatives for the conditioned finite comparison diagrams
specified in §8.

The ordered coefficient bound, convergent record flows and Einstein
realization were established as part of the projective twistor
correspondence in Isthmus (2026, Lemma 3.2, Theorem 3.3, Proposition
4.1, Appendix A and Theorem B.4). The flow proof there also gives
holomorphic maps into smaller-radius sequence spaces. This paper
develops the local realization theorem in a self-contained form. We
expand the normalized diagram estimates and formulate the record-flow
theorem on the analytic scale. We construct the metric directly from
radial curvature by the Jacobi equation, and prove that the normal
metric and its inverse depend holomorphically on the records in Banach
spaces of bounded holomorphic tensors. The point and full-spin
comparisons and Lorentz real form correspond to the metric statements of
Isthmus (2026, secs. 8.3--8.4).

Let \(S,S'\) be marked complex symplectic two-spaces, \(V=S\otimes S'\)
and \(\eta\) the induced tangent form. The free record spaces are \[
 W_r^+=\operatorname{Sym}^{r+4}S^*\otimes\operatorname{Sym}^{r}S'^*,
 \qquad W_r^-=\operatorname{Sym}^{r}S^*\otimes\operatorname{Sym}^{r+4}S'^*.
\] For an Einstein metric these are the separately symmetric parts of
\(\nabla^r\Psi^\pm\) at the marked point. An analytic abstract record
consists of independent \(U_r^\pm\in W_r^\pm\) satisfying
\(\|U_r^\pm\|\le Cr!R_0^{-r}\) for some \(C,R_0>0\), together with a
numerical constant \(\Lambda\). Equivalently it belongs to one of the
Banach spaces \(\mathcal B_R\) defined in §2.

The main results are summarized below and proved in Theorem
\hyperref[thm:es-realization]{7.2}, Proposition
\hyperref[prop:es-uniform-realization]{7.3} and §8.

\textbf{Theorem A.} Every analytic abstract record
\(X=(U^+,U^-,\Lambda)\) is realized by a marked holomorphic Einstein
four-germ with \(\operatorname{Ric}(g)=\Lambda g\). Its entire free
curvature record is the prescribed one, and it is unique up to the
marked isometry. Every marked holomorphic Einstein four-germ arises in
this way. For each \(R,M>0\) there is a common normal-coordinate
neighborhood on the record ball \(\|X\|_R\le M\), with common bounds for
the metric and its inverse. The normal realization, its coordinate maps
and radial spin frames depend holomorphically on the records, with the
spatial restrictions specified in Proposition
\hyperref[prop:es-uniform-realization]{7.3}.

\textbf{Theorem B.} Short changes of point and full spin frame are
represented by the convergent record flows. They give exactly the local
isometries with their specified spin lifts, and their finite conditioned
composition diagrams have common representatives on bounded record
families. The Lorentz-real records are precisely those of real-analytic
Lorentzian Einstein four-germs.

The conditioning in §8 bounds the spin maps and their inverses and keeps
the coordinate maps and specified composites inside their assigned
normal neighborhoods with uniformly positive buffer and inverse margins.

We first complete the free records to all ordered curvature derivatives.
We verify the polynomial identities on analytic Einstein metrics
realizing finite record prefixes. This formal completion defines the
horizontal derivations corresponding to a change of base point: in a
direction \(a\) they act by \[
 \mathsf H_aU_r=i_aU_{r+1}+Q_{r,a}(U,\Lambda).
\] We use auxiliary Hermitian tensor norms to estimate the derivations
on an analytic record scale. The leading shift loses radius, while the
nonlinear terms come from the curvature commutators in the exact-set
recursion. To integrate these derivations, we need estimates uniform in
derivative order and curvature degree. Lemma
\hyperref[lem:es-transfer]{5.4} supplies an inverse-order factor
\(1/(t+1)\) whenever an order-\(t\) correction is inserted into the
recursion. A common contraction-diagram mass controls these insertions,
and a rescaling absorbs the summable linear transfer loss. The
inverse-order gain then closes the curvature-degree induction. The
resulting ordered degree-\(d\) coefficients satisfy \[
 \|B^{[d]}_{r,\mathbf j}\|
 \le ab^{d-2}r\,\frac{r!}{\prod_i j_i!},
 \qquad j_1+\cdots+j_d=r+3-2d,
\] for structural constants \(a,b>0\). A scalar majorant equation gives
normally convergent horizontal flows on smaller positive record radii,
with holomorphic initial-value dependence.

For a small tangent vector \(x\), the flow \(\Phi_x^t\) of
\(\mathsf H_x\) gives the analytic radial curvature
\(\mathcal C_0(\Phi_x^tX)\). The initialized radial Jacobi equation
constructs the metric in marked normal coordinates. Its Taylor series
agrees with the formal normal metric determined by the completed
curvature list, so the formal Einstein and curvature-readout identities
become identities of this analytic metric. The uniform flow and Jacobi
estimates give the common normal domains and Banach-holomorphic
realization of Proposition \hyperref[prop:es-uniform-realization]{7.3}.

Penrose introduced exact sets and their constraint-free null-cone
interpretation for interacting fields, including the Einstein equations
(Penrose 1980, secs. 2--3). The spinor coordinates of Anderson and Torre
describe the free components of prolonged Ricci-flat jets and the
leading part of their total derivatives (Anderson and Torre 1996, sec.
2C, Theorems 2.7--2.8 and Proposition 2.10). Torre explicitly
distinguishes that formal parametrization from convergence of the
resulting metric series (Torre 1995, sec. 2).

Related realization problems prescribe algebraic curvature at a point,
compatible lists of curvature derivatives, or an analytic curvature map.
Gasqui proves formal integrability and involutivity of the Einstein
system and realizes prescribed algebraic Einstein curvature at a point,
with the extension to arbitrary signatures stated after his Theorem 4.2
(Gasqui 1982, sec. 4). Kowalski and Belger construct analytic metrics in
normal coordinates from compatible sequences of all covariant curvature
derivatives satisfying growth inequalities (Kowalski and Belger 1994).
The classical normal-jet correspondence expresses the normal metric jet
by universal curvature polynomials; a coordinate-free formulation is
given by Jentsch (2026, Theorems 1 and 3). Basurto-Arzate realizes an
analytic curvature map satisfying a radial structure consistency
identity by a torsion-free connection (Basurto-Arzate 2026, Theorem 1).
Here the input is the two free chiral sequences: completion and
convergent record flows supply the analytic radial curvature used in the
Jacobi construction.

Friedrich proves a factorial null-data criterion at spatial infinity for
static vacuum solutions with a fixed nonzero ADM mass (Friedrich 2007,
Theorem 1.1 and Lemma 3.2). Aceña extends this reconstruction to
asymptotically flat stationary vacuum solutions with nonzero mass
monopole (Aceña 2009, Theorem 1.1). His two null-data sequences are the
symmetric trace-free derivatives of a conformally rescaled scalar
potential \(\phi_S\) and of the trace-free Ricci tensor of the
three-dimensional conformal quotient metric. Aceña subsequently obtains
a null-data characterization for three-dimensional negative-definite
geometry coupled to scalar fields satisfying the coefficient conditions
of his system (Aceña 2010, Theorem 1.1 and §7).

In four dimensions, Friedrich identifies the corresponding symmetric
vertex spinors as null data for the rescaled conformal Weyl field,
determines the formal conformal-vacuum jets from them, and constructs
smooth fields satisfying the equations to all orders at the vertex
(Friedrich 2013, sec. 5.1, Lemma 6.1 and Proposition 8.1). He also
remarks in §1 that the analysis simplifies for a finite cone with a
regular vacuum vertex. Chruściel and Paetz prove existence and
uniqueness of smooth vacuum space-times from freely prescribed
cone-smooth radiation data at past null infinity (Chruściel and Paetz
2013, Theorem 1.1). The present theorem realizes two independent chiral
Weyl records at a regular point by a holomorphic Einstein metric with
arbitrary numerical \(\Lambda\), with factorial record estimates and
uniform holomorphic dependence.

Sections 2--4 establish the record spaces, universal completion and
normalized horizontal recursion. Section 5 proves the coefficient
estimate, §6 integrates the record fields, and §7 constructs the normal
Einstein metric by the radial Jacobi equation. Section 8 gives the point
and spin comparisons and the Lorentz real form. Appendix A proves
finite-prefix realizability, and Appendix B records low-order forms of
the recursion.

\protect\phantomsection\label{sec:es-records}

\subsection{2. Curvature records and analytic
scales}\label{curvature-records-and-analytic-scales}

Let \(S,S'\) be complex symplectic two-spaces with fixed ordered bases.
Put \(V=S\otimes S'\) and
\(\eta_{AA'BB'}=\epsilon_{AB}\epsilon_{A'B'}\). Our index conventions
are
\begin{equation}\protect\phantomsection\label{eq:es-spin-conventions}{\begin{gathered}
\epsilon_{AB}=\epsilon_{A'B'}=
 \begin{pmatrix}0&1\\-1&0\end{pmatrix}, \\
\epsilon^{AB}=\epsilon^{A'B'}=
 \begin{pmatrix}0&-1\\1&0\end{pmatrix}, \\
v^A=\epsilon^{AB}v_B,\quad v_A=\epsilon_{AB}v^B .
\end{gathered}}\end{equation} The primed operations have the same order. A marked
holomorphic metric germ consists of a metric near a point \(o\) and a
full spin marking: an identification of \(S,S'\) with the spin spaces at
\(o\), inducing the tangent isometry \(V\to T_oM\). A marked isometry
sends the base point to the base point and intertwines these markings by
its specified spin lift. Germs are identified by restriction to smaller
neighborhoods. Spin lifts of comparisons are retained as part of their
data. Write
\(\operatorname{Spin}(V,\eta)=\operatorname{SL}(S)\times\operatorname{SL}(S')\).

Our curvature operator is \(R(a,b)=[\nabla_a,\nabla_b]-\nabla_{[a,b]}\),
and \(\operatorname{Ric}(b,c)=\operatorname{tr}(a\mapsto R(a,b)c)\).
Lowering the output in the last slot gives
\(R_{abcd}=g(R(e_a,e_b)e_c,e_d)\). With
\(P=\tfrac12(\operatorname{Ric}-\operatorname{Scal}\,g/6)\), define
\begin{equation}\protect\phantomsection\label{eq:es-curvature-decomposition}{
 \begin{aligned}
 C_{abcd}&=R_{abcd}-g_{ad}P_{bc}+g_{bd}P_{ac}
                         +g_{ac}P_{bd}-g_{bc}P_{ad},\\
 C_{AA'BB'CC'DD'}
 &=-\Psi^+_{ABCD}\epsilon_{A'B'}\epsilon_{C'D'}
   -\Psi^-_{A'B'C'D'}\epsilon_{AB}\epsilon_{CD} .
 \end{aligned}
}\end{equation} The two spinors are therefore recovered by \[
 \begin{aligned}
 \Psi^+_{ABCD}&=-\tfrac14\epsilon^{A'B'}\epsilon^{C'D'}C_{AA'BB'CC'DD'},\\
 \Psi^-_{A'B'C'D'}&=-\tfrac14\epsilon^{AB}\epsilon^{CD}C_{AA'BB'CC'DD'} .
 \end{aligned}
\] Throughout, \(\Lambda\) is the coefficient in
\(\operatorname{Ric}(g)=\Lambda g\).

\pandocstatementheading{\protect\phantomsection\label{def:es-records}Definition 2.1.}

For \(r\ge0\) let
\begin{equation}\protect\phantomsection\label{eq:es-record-spaces}{
 W_r^+=\operatorname{Sym}^{r+4}S^*\otimes\operatorname{Sym}^{r}S'^*,
 \qquad
 W_r^-=\operatorname{Sym}^{r}S^*\otimes\operatorname{Sym}^{r+4}S'^* .
}\end{equation} The free records of a marked Einstein metric are
\begin{equation}\protect\phantomsection\label{eq:es-free-records}{
 U_r^+=\operatorname{Sym}_{A_1\cdots A_{r+4}}
       \operatorname{Sym}_{A'_1\cdots A'_r}
       \nabla_{A_1A'_1}\cdots\nabla_{A_rA'_r}
       \Psi^+_{A_{r+1}\cdots A_{r+4}}\big|_o,
}\end{equation} and the corresponding negative expression with the two
spin factors exchanged. Symmetrization is averaging. An abstract record
is a pair of arbitrary sequences \(U_r^\pm\in W_r^\pm\) and a scalar
\(\Lambda\). A finite prefix retains \(0\le r\le N\).

Give both spin spaces the auxiliary Hermitian norms for which their
fixed bases are orthonormal. Symmetric tensors carry the subspace norms
of full tensor powers; direct sums carry the orthogonal direct-sum norm.
These norms are used for analytic estimates. Put
\(X_0=(U_0^+,U_0^-,\Lambda)\) and \(X_r=(U_r^+,U_r^-)\) for \(r>0\).

\pandocstatementheading{\protect\phantomsection\label{def:es-record-scale}Definition 2.2.}

For \(R>0\) the analytic record space \(\mathcal B_R\) is the Banach
space with norm
\begin{equation}\protect\phantomsection\label{eq:es-record-norm}{
 \|X\|_R=\sum_{r\ge0}\frac{R^r}{r!}\|X_r\|.
}\end{equation} For \(R>s>0\), write
\(\iota_{R,s}:\mathcal B_R\to\mathcal B_s\) for the continuous
inclusion. A bounded analytic record family has a common \(R>0\) and a
common bound for this norm. A holomorphic family is a holomorphic map
into such a space on a smaller parameter neighborhood. For a fixed
numerical constant one takes the corresponding affine slice of
\(\mathcal B_R\).

\pandocstatementheading{\protect\phantomsection\label{lem:es-record-growth}Lemma 2.3.}

An abstract record belongs to some \(\mathcal B_R\) if and only if there
are \(C,R_0>0\) such that \(\|U_r^\pm\|_{\mathrm{ev}}\le Cr!R_0^{-r}\),
where
\begin{equation}\protect\phantomsection\label{eq:es-evaluation-norm}{
 \|T\|_{\mathrm{ev}}=
 \sup_{|\xi|,|\pi|\le1}|T(\xi^p,\pi^q)|,
 \qquad T\in\operatorname{Sym}^pS^*\otimes\operatorname{Sym}^qS'^* .
}\end{equation} The two descriptions of bounded families agree after a
fixed radius decrease.

\emph{Proof.} Coherent evaluation and the unit-sphere integral give
\begin{equation}\protect\phantomsection\label{eq:es-coherent-norm}{
 \|T\|_{\mathrm{ev}}\le\|T\|
 \le\sqrt{(p+1)(q+1)}\,\|T\|_{\mathrm{ev}}.
}\end{equation} With probability surface measure on each unit spin
sphere, phase orthogonality and
\(\int_0^1t^j(1-t)^{p-j}\,\mathrm dt=j!(p-j)!/(p+1)!\) show that the
squared coherent integral is \(\|T\|^2/((p+1)(q+1))\). The first
inequality is Cauchy--Schwarz. For \(W_r^\pm\) the square-root factor is
at most \(r+3\). At \(R=R_0/2\) the two series therefore have sum at
most \(16C\), and \(\|X\|_R\le16C+|\Lambda|\). Conversely each summand
in the record norm is bounded by the norm itself. ∎

For the homogeneity count in the curvature completion, assign scaling
weight \(r+2\) to \(U_r^\pm\) and weight two to \(\Lambda\). Curvature
degree counts the record or scalar factors in a monomial; its input
positions remain ordered and retain their spin types.

\pandocstatementheading{\protect\phantomsection\label{def:es-lorentz-involution}Definition 2.4.}

Choose the Lorentz spin identification for which real vectors in \(V\)
are Hermitian matrices and \(\eta(v,v)=2\det(v^{AA'})\), of signature
\((+,-,-,-)\). Let \(\mathfrak c\) conjugate coefficients, exchange the
spin factors and the two chiral records, and conjugate \(\Lambda\). A
Lorentz-real record is fixed by \(\mathfrak c\); in the fixed bases this
says \(U^-_r=\overline{U^+_r}\) and \(\Lambda\in\mathbb R\).

\protect\phantomsection\label{sec:es-completion}

\subsection{3. Highest compatibility and curvature
completion}\label{highest-compatibility-and-curvature-completion}

The algebraic and differential Bianchi identities, together with the
commutator of covariant derivatives, determine the part of each ordered
curvature derivative complementary to the free records. We first
identify the homogeneous system that occurs when the lower derivatives
have been fixed. Einstein contractions are the differentiated equations
\(\nabla^k(\operatorname{Ric}-\Lambda g)=0\), \(k\ge0\), with
\(\Lambda\) constant.

\pandocstatementheading{\protect\phantomsection\label{lem:es-highest-kernel}Lemma 3.1.}

Fix the curvature derivatives of orders less than \(r\), with the
Einstein Ricci tensor \(\Lambda\eta\). The homogeneous difference of two
order-\(r\) solutions of the differentiated curvature symmetries,
Bianchi identities, Einstein contractions and derivative-commutator
identities is precisely \(W_r^+\oplus W_r^-\). The separately symmetric
free projection restricts to the identity on this kernel.

\emph{Proof.} At order zero the curvature decomposition
eq.~\ref{eq:es-curvature-decomposition} gives the assertion. Suppose
\(r\ge1\) and let \(T\) be the positive Weyl component of a highest
difference. Interchanging adjacent derivative pairs changes an ordered
derivative by curvature acting on a lower derivative, so those pairs
commute in \(T\). The Ricci derivatives in the difference vanish. The
differentiated Einstein Bianchi identity therefore makes the alternating
contraction between any derivative unprimed index and any of the four
Weyl unprimed indices zero.

For two specified slots in a two-dimensional spin space, the convention
eq.~\ref{eq:es-spin-conventions} gives
\begin{equation}\protect\phantomsection\label{eq:es-two-slot-symmetry}{
 T_{AB\cdots}-T_{BA\cdots}
   =-\epsilon_{AB}\epsilon^{CD}T_{CD\cdots}.
}\end{equation} Thus the vanishing contractions make \(T\) invariant
under interchanging each derivative unprimed index with a Weyl unprimed
index. Those interchanges, together with the symmetry of the four Weyl
indices, generate every permutation of the \(r+4\) unprimed slots. Hence
all unprimed slots are symmetric. Since the derivative pairs are
symmetric, exchanging two derivative primed slots can now be compensated
by exchanging their already symmetric unprimed slots. All \(r\) primed
slots are consequently symmetric as well. The positive difference is an
element of
\(\operatorname{Sym}^{r+4}S^*\otimes\operatorname{Sym}^rS'^*\).

The negative Weyl component has the factors exchanged and belongs to
\(W_r^-\). Conversely, a tensor in either of these separately symmetric
spaces has zero alternating contractions and commuting derivative pairs,
and satisfies every homogeneous highest-order identity. Its free
projection is itself. ∎

At each finite order, these identities form a system between fixed
finite-dimensional tensor spaces. Its coefficients use the fixed
alternating tensors and tangent form; products of lower derivatives
appear in the inhomogeneous terms through the derivative-commutator
identities.

\pandocstatementheading{\protect\phantomsection\label{prop:es-completion}Proposition 3.2.}

For every \(r\ge0\) there is a unique universal polynomial tensor
\(\mathcal C_r(U,\Lambda)\) that completes the two free records through
order \(r\) to the ordered Einstein curvature derivative of order \(r\).
These polynomials satisfy all differentiated curvature symmetries,
Bianchi identities, Einstein contractions and derivative-commutator
identities. They are equivariant for the two complex spin groups and
commute with the Lorentz coefficient involution.

\emph{Proof.} At order zero use the Weyl decomposition and the
prescribed Ricci part. Inductively, collect the order-\(r\) identities
and the specified free projection into a finite constant linear system
\begin{equation}\protect\phantomsection\label{eq:es-finite-completion-system}{
 A_r C_r=b_r(U_0,\ldots,U_r,\Lambda).
}\end{equation} The already constructed lower-order polynomial tensors
make the right-hand side polynomial. Lemma 3.1 makes \(A_r\) injective,
so a constant left inverse defines a candidate polynomial
\(\mathcal C_r\).

Lemma \hyperref[lem:es-finite-witness]{A.1}, proved by finite analytic
Cauchy systems and a finite-dimensional inverse theorem, realizes every
finite Lorentz-real prefix and every real scalar by an analytic Einstein
metric. On that metric, the actual order-\(r\) derivative solves the
full system and has the prescribed free projection. Injectivity makes it
equal to the candidate. Hence the residual \(A_r\mathcal C_r-b_r\)
vanishes on the real form \(U^-=\overline{U^+}\),
\(\Lambda\in\mathbb R\).

Choose paired bases for the two chiral spaces and write
\(U^+=x+\mathrm i y\), \(U^-=x-\mathrm i y\). A complex polynomial in
\((x,y,\Lambda)\) that vanishes for all real values of these variables
is identically zero. The residual therefore vanishes for independent
complex records and complex \(\Lambda\). The candidate satisfies the
full system, and uniqueness makes it independent of the chosen left
inverse.

Spin transformations preserve the defining identities and transform the
free projection, so uniqueness gives equivariance of the completion.
Conjugation with exchange of the two spin factors preserves the same
defining equations and gives the Lorentz assertion. ∎

For a covariant tensor \(T\) on \(V\), write \(\mathsf V_BT\) for
insertion of \(B\in\mathfrak{so}(V,\eta)\) into each argument, summed
over all arguments. Define \(\mathsf H_a\) on the free coordinates by
taking the free projection of the next completed derivative with first
derivative argument \(a\in V\). Write \(i_aU_{r+1}\) for contraction of
one derivative pair of the separately symmetric record with \(a\). These
definitions extend by the polynomial product rule to the algebra of
cylinder polynomials in the records, that is, polynomials depending on
finitely many components. They fix \(\Lambda\).

\pandocstatementheading{\protect\phantomsection\label{prop:es-polynomial-derivations}Proposition 3.3.}

The resulting polynomial derivations satisfy
\begin{equation}\protect\phantomsection\label{eq:es-cartan-identities}{
 \begin{aligned}
 [\mathsf H_a,\mathsf H_b]&=-\mathsf V_{R_U(a,b)},&
 [\mathsf V_B,\mathsf H_a]&=\mathsf H_{Ba},&
 [\mathsf V_B,\mathsf V_C]&=\mathsf V_{[B,C]},\\
 \mathsf H_a\Lambda&=\mathsf V_B\Lambda=0,&
 \mathsf H_a\mathcal C_r&=\mathcal C_{r+1}(a;\,\cdot),
 \end{aligned}
}\end{equation} where
\begin{equation}\protect\phantomsection\label{eq:es-algebraic-curvature}{
 R_U(a,b)c=W_U(a,b)c+
 \frac{\Lambda}{3}\bigl(a\eta(b,c)-b\eta(a,c)\bigr),
}\end{equation} and \(W_U\) is the Weyl tensor with chiral spinors
\(U_0^+,U_0^-\). Moreover,
\begin{equation}\protect\phantomsection\label{eq:es-horizontal-polynomial}{
 \mathsf H_aU_r=i_aU_{r+1}+Q_{r,a}(U,\Lambda),\qquad Q_0=0,
}\end{equation} where \(Q_r\) depends only on
\(U_0,\ldots,U_{r-1},\Lambda\). At curvature degree \(d\ge2\), each
ordered monomial of \(Q_r\) has input record orders satisfying
\begin{equation}\protect\phantomsection\label{eq:es-completion-weight}{
 j_1+\cdots+j_d=r+3-2d;
}\end{equation} a scalar \(\Lambda\) occupies an input position of
derivative order zero.

\emph{Proof.} Each component of each displayed identity depends on a
finite prefix. Realize a prefix long enough to include all derivatives
appearing in that component by Lemma
\hyperref[lem:es-finite-witness]{A.1}. Proposition 3.2 applies to the
records at every point and frame of this witness. On its spin frame
bundle, \(\mathsf H_a\) is therefore the horizontal component readout
and \(\mathsf V_B\) is the frame readout. The horizontal commutator acts
on every covariant slot by minus curvature insertion. The
horizontal--vertical and vertical--vertical brackets give the other two
commutators with the stated pullback convention. Differentiating an
ordered curvature derivative gives the last identity, including its
derivative slots. Thus each component identity holds on every finite
Lorentz-real prefix. The polynomial continuation used in Proposition 3.2
proves all of them on independent complex records.

Under the dilation \(g_s=s^{-2}\delta_s^*g\), \(\delta_s(x)=sx\), the
tangent metric at the mark is fixed, \(U_j\) has weight \(j+2\) and
\(\Lambda\) has weight two. The value of \(\mathsf H_aU_r\) has weight
\(r+3\). Its linear term is \(i_aU_{r+1}\): in the homogeneous linear
system of Lemma 3.1, the highest derivative is the separately symmetric
tensor itself, and its free projection is an average. Every remaining
monomial has degree \(d\ge2\) and therefore has weight
\(j_1+\cdots+j_d+2d\). This gives the degree relation. It implies
\(j_i\le r-1\) for every nonlinear input, and gives \(Q_0=0\). ∎

\protect\phantomsection\label{sec:es-derivations}

\subsection{4. The Einstein Cartan
derivations}\label{the-einstein-cartan-derivations}

Proposition \hyperref[prop:es-polynomial-derivations]{3.3} defines the
horizontal derivations \(\mathsf H_a\) and vertical derivations
\(\mathsf V_B\) on cylinder polynomials. The horizontal field has the
form \(\mathsf H_aU_r=i_aU_{r+1}+Q_{r,a}\). Resolving the antisymmetric
part of one additional covariant derivative gives the nonlinear
correction.

Use the full-tensor Hermitian norms of §2, with the fixed alternating
spin bases orthonormal. Symmetrizations are averaging projections. All
adjoints below are adjoints of constant finite matrices; they do not
introduce conjugated record variables.

Fix either chirality and write \(W_r\) for its record space. For
\(r\ge1\), regrouping one spinor of each kind gives an isometric
inclusion \(j_r:W_r\to V^*\otimes W_{r-1}\). Let \(J_r=1\otimes j_r\)
and let \(\tau\) exchange the two outer derivative pairs after that
inclusion. Set \(\delta_r=(1-\tau)J_r\), without dividing by two. For
the positive chirality write \(n=r+4,m=r\); exchange them for the
negative one. The upper and lower projections in the two spin factors
are denoted by \(P_{++},P_{-+},P_{+-},P_{--}\). In each spin factor, the
upper summand is fully symmetric and the lower summand is its
complement; §5.1 gives their polynomial normalizations.

\pandocstatementheading{\protect\phantomsection\label{lem:es-spencer}Lemma 4.1.}

For the original derivation \(\mathsf H_aU_r=i_aU_{r+1}+Q_{r,a}\),
\(r\ge0\), the initial case is \(Q_0=0\). For \(r\ge1\) one has the
following identities, with the last sum interpreted as zero when
\(r=1\): \begin{equation}\protect\phantomsection\label{eq:es-recursion}{
 \begin{aligned}
 \ker\delta_r&=\operatorname{im}P_{++},\\
 M_r&=\frac n{n+1}P_{-+}+\frac m{m+1}P_{+-}
                         +\frac{nm}{nm-1}P_{--},\\
 Q_r&=M_rJ_r^*(\tau\,\mathsf H Q_{r-1}-K_r),\\
 K_r(a,b)&=\mathsf V_{R_U(a,b)}U_{r-1},\\
 (\mathsf H Q_{r-1})(a,b)
  &=\sum_{j=0}^{r-2}D_{U_j}Q_{r-1,b}
                         [i_aU_{j+1}+Q_{j,a}].
 \end{aligned}
}\end{equation} \emph{Proof.} On \(A\otimes\operatorname{Sym}^nA\), let
\(P_n\) be the orthogonal full-symmetrization projection onto
\(\operatorname{Sym}^{n+1}A\). Compression of an exchange of the outer
slot with an inner slot is \(C_n=((n+1)P_n-1)/n\). Indeed full
symmetrization is \((1+\sum_{i=1}^n\tau_{0i})/(n+1)\), and the
compressed inner exchanges agree. Its two eigenvalues are \(1,-1/n\).
Thus
\begin{equation}\protect\phantomsection\label{eq:es-spencer-spectrum}{
 \delta_r^*\delta_r=2(1-C_n\otimes C_m)
}\end{equation} has nonzero eigenvalues \(2(1+1/n)\), \(2(1+1/m)\) and
\(2(1-1/(nm))\). Their minimum is at least \(8/5\). On antisymmetric
sources the complementary inverse of \(\delta_r\) is \(M_rJ_r^*\).

Apply the actual commutator in Proposition
\hyperref[prop:es-polynomial-derivations]{3.3} to \(U_{r-1}\). The
completely upper \(U_{r+1}\) term cancels and gives
\(\delta_rQ_r=-K_r-(1-\tau)\mathsf H Q_{r-1}\). Differentiating the
constant zero-upper-projection identity for \(Q_{r-1}\) gives
\(J_r^*\mathsf H Q_{r-1}=0\). The commutator identity places the
antisymmetric source in the image of \(\delta_r\). Applying its
complementary inverse recovers the already defined \(Q_r\) and proves
the formula. Here \((\mathsf H Q_{r-1})(a,b)=\mathsf H_a(Q_{r-1,b})\)
has direction slot \(a\) first and the coefficient's own slot \(b\)
second; \(J_r^*\tau\) retains \(b\) and symmetrizes \(a\) into the
output record. The last line is the ordinary polynomial product rule;
\(\mathsf H_a\Lambda=0\). The curvature in \(K_r\) acts on all \(2r+2\)
spinor slots, including the derivative slots. ∎

Appendix B displays the first three corrections and the first possible
insertion of one nonlinear correction into another.

\protect\phantomsection\label{sec:es-coefficients}

\subsection{5. Uniform nonlinear coefficient
estimates}\label{uniform-nonlinear-coefficient-estimates}

The convergence argument requires coefficient bounds uniform in
derivative order and curvature degree. We obtain them by measuring
contraction diagrams in one mass and retaining the inverse-order gain in
nonlinear insertions. Throughout this section, symmetric tensors have
the full-tensor Hermitian norms of §2, and every symmetrization is an
average. Structural constants depend only on the two fixed spin spaces
and the curvature convention.

At curvature degree \(d\), retain the order and chirality of each input
position in the polynomial expansion of \(Q_r\). A scalar \(\Lambda\)
position has order zero. The dilation weights in Proposition
\hyperref[prop:es-polynomial-derivations]{3.3} give
\begin{equation}\protect\phantomsection\label{eq:es-ordered-expansion}{
 Q_{r,a}(X)=\sum_{d\ge2}\ 
 \sum_{|\mathbf j|=r+3-2d}
 B^{[d]}_{r,\mathbf j}(a;X_{j_1},\ldots,X_{j_d}).
}\end{equation} We first estimate coefficients with fixed input and
output chiralities; they are grouped at the end of §5.3. Recursion
(\ref{eq:es-recursion}) fixes the ordered expansion: differentiation
replaces an input at its existing position, and an insertion preserves
the order of its own inputs. This convention determines coefficients of
the commutative polynomial without a symmetric polarization or a factor
\(1/d!\).

\protect\phantomsection\label{sec:es-contractions}

\subsubsection{5.1. Normalized
contractions}\label{normalized-contractions}

In one spin factor, express a symmetric tensor by its polynomial on a
vector \(u=(u^0,u^1)\). For a degree-\(q\) polynomial with an additional
covariant index, write it as \(T_a(u)\). Put
\begin{equation}\protect\phantomsection\label{eq:es-upper-inclusion}{
 (\operatorname{up}_{q+1}g)_a=\frac{\partial_ag}{q+1},\qquad
 g=u^aT_a.
}\end{equation} For \(q\ge1\) also put
\begin{equation}\protect\phantomsection\label{eq:es-lower-inclusion}{
 (\iota_qf)_0=-u^1f,\qquad
 (\iota_qf)_1=u^0f,\qquad
 \operatorname{strip}_qT
 =\frac{\partial_0T_1-\partial_1T_0}{q+1}.
}\end{equation} Then
\(T=\operatorname{up}_{q+1}g+\iota_q\operatorname{strip}_qT\). For
\(q=0\) only the upper summand occurs. The upper inclusion is the
isometric inclusion of a completely symmetric tensor in the full tensor
product. For \(q\ge1\), the squared norm of the lower inclusion is
\((q+1)/q\le2\): on the unit highest vector \((u^0)^{q-1}\) its two
components have squared norms \(1/q\) and \(1\), and unitary
equivariance gives the general formula.

A bare contraction diagram, or atom, has a fixed ordered list of
symmetric inputs with their two spin valences. In each spin factor,
choose a pairing of some input slots, attach the alternating tensor
\(\epsilon^{ab}\) to every pair with a specified orientation, and pass
every remaining slot to an output of the same spin type. Average the
output slots separately in the two factors. An input of scalar type
\(\Lambda\) has no spin slots. Each alternating tensor joins two input
slots directly; there are no chains of alternating tensors or closed
components.

Slot selection is normalized: choosing \(h\) slots of a degree-\(p\)
symmetric polynomial means applying
\(\partial_{a_1}\cdots\partial_{a_h}/(p)_h\), where
\((p)_h=p(p-1)\cdots(p-h+1)\). Thus labeled slot choices represent the
same averaged tensor operation. Reversing the orientation of a
contracted pair changes its sign. A contraction between two slots of the
same symmetric input is zero. An atom with \(c\) alternating
contractions has operator norm at most \(2^{c/2}\), since each
alternating tensor has full-tensor norm \(\sqrt2\) and slot selections
and symmetrizations have norm at most one.

An output block is restored by applying \(\operatorname{up}\) in its
upper factors and \(\iota\) in its lower factors. To compose restored
diagrams, resolve each symmetrization as its average of slot
permutations, and glue the matching slots. Identity tensors
\(\delta^a_b\) and alternating tensors then have two incident
half-edges, while the uncontracted input and output slots are terminals.
A \emph{wire} is a connected component of this graph after the glued
internal slots are suppressed; it is a path or a circle. An \emph{opened
wire} at a lower insertion is obtained when a selected derivative hits
its raw factor \(\epsilon_{bc}u^b\): compose this alternating tensor
with the incident contracted pair, replace their composition by the
signed identity, and pass the partner slot to the symmetric output.
Lemma 5.1 records the normalization of this operation.

\pandocstatementheading{\protect\phantomsection\label{lem:es-opening}Lemma 5.1.}

Replace a degree-\(p\) input of a normalized contraction diagram by
\((\iota_pf)_c\), and suppose that \(h\) of its slots are contracted.
Let \(D_c(u)\) be the resulting degree-\(N\) polynomial in the free
output variable of this spin factor. If \(E_j\) is the diagram obtained
by opening the \(j\)th incident contraction, with its orientation sign
\(\varepsilon_j\), then
\begin{equation}\protect\phantomsection\label{eq:es-opened-edge}{
 u^cD_c=\frac1p\sum_{j=1}^h\varepsilon_jE_j.
}\end{equation} For \(N\ge1\) this also gives
\begin{equation}\protect\phantomsection\label{eq:es-opened-edge-derivative}{
 \frac{u^c\partial_aD_c}{N}
 =\frac{N+1}{N}(\operatorname{up}_{N+1})_a
       \left(\frac1p\sum_j\varepsilon_jE_j\right)
       -\frac{D_a}{N}.
}\end{equation} For a linear shift, in which the input is replaced by
\(\operatorname{up}_{p+1}g\), contraction of the new index with \(u\)
has coefficient one. An old lower output has normalized coefficient
bound one. For an old upper output, the lower leakage obtained by
stripping has coefficients of absolute value
\begin{equation}\protect\phantomsection\label{eq:es-linear-leakage-coefficient}{
 \frac{p_j-h_j}{N(N+1)}
}\end{equation} for the other inputs \(j\), where \(p_j-h_j\) is their
number of free slots. The sum of these coefficients is at most
\(1/(N+1)\).

\emph{Proof.} In the fixed alternating convention,
\((\iota_pf)_c=\epsilon_{bc}u^bf\). The normalized product rule reads
\begin{equation}\protect\phantomsection\label{eq:es-normalized-product}{
 \frac{\partial_{a_1}\cdots\partial_{a_h}
                  (\epsilon_{bc}u^bf)}{(p)_h}
 =\frac{p-h}{p}\epsilon_{bc}u^b
       \frac{\partial_{a_1}\cdots\partial_{a_h}f}{(p-1)_h}
  +\frac1p\sum_{j=1}^h\epsilon_{a_jc}
       \frac{\prod_{\ell\ne j}\partial_{a_\ell}f}{(p-1)_{h-1}}.
}\end{equation} The first term is omitted when \(h=p\), and the sum is
empty when \(h=0\). Multiplication by \(u^c\) kills the first term. In
each remaining term the two alternating tensors at the opened
contraction compose to a signed identity. If the partner input has
degree \(p_j\) and \(h_j\) selected slots, the released derivative gives
the factor \(p_j-h_j+1\). Its normalization is
\begin{equation}\protect\phantomsection\label{eq:es-partner-normalization}{
 \frac{p_j-h_j+1}{(p_j)_{h_j}}
       =\frac1{(p_j)_{h_j-1}}.
}\end{equation} This leaves precisely the normalized selection with
\(h_j-1\) slots. Thus the coefficient of the opened diagram is exactly
\(\varepsilon_j/p\), without a partner-valence factor. Differentiating
\(u^cD_c\) proves eq.~\ref{eq:es-opened-edge-derivative}.

For the upper replacement, the selected polynomial derivatives are
\(\partial_{a_1}\cdots\partial_{a_h}\partial_cg/((p+1)(p)_h)\). Euler's
identity after contraction with \(u^c\) gives the multiplier
\begin{equation}\protect\phantomsection\label{eq:es-unit-linear-shift}{
 \frac{(p+1-h)(p+1)_h}{(p+1)(p)_h}=1.
}\end{equation} To compute leakage from an old upper output, write its
restored outer index as a normalized derivative. Stripping the
transferred output antisymmetrizes this derivative with the new
derivative index. When both derivatives land on the distinguished
shifted input, their contribution vanishes. A derivative landing on
another input \(j\) can select any of its \(p_j-h_j\) remaining free
slots. The two output normalizations give the denominator \(N(N+1)\) in
eq.~\ref{eq:es-linear-leakage-coefficient}. The remaining free-slot
counts sum to at most \(N\), proving the total bound. Raw lower output
factors are independent of the distinguished input, so an old lower
output stays lower with unit bound. These arguments are tensor-valued in
the other spin factor. ∎

For example, a single bare contraction of degree-\(p\) and degree-\(q\)
inputs is \(\epsilon^{ab}(\partial_aP)(\partial_bQ)/(pq)\). Replacing
the first input by \(\epsilon_{dc}u^df\) gives, in the alternating
convention of §2,
\begin{equation}\protect\phantomsection\label{eq:es-single-opened-edge-example}{
 u^c\frac{\epsilon^{ab}}{pq}
       \partial_a(\epsilon_{dc}u^df)\,\partial_bQ
 =-\frac{u^b\partial_bQ}{pq}f
 =-\frac1p fQ.
}\end{equation} The factor \(q\) from Euler's identity cancels the
partner input's normalization. This is the cancellation used at every
opened contraction in Lemma 5.1.

The one-derivative space has four upper/lower blocks. The \(++\) block
of \(Q_r\) is zero by its definition. For the positive chirality, strip
the raw lower factors and retain the upper potentials in the other
factors. The three remaining output blocks have valences
\begin{equation}\protect\phantomsection\label{eq:es-stripped-valences}{
 \begin{array}{c|ccc}
 \nu&A=-+&B=+-&C=--\\ \hline
 \text{stripped valence}&(r+3,r+1)&(r+5,r-1)&(r+3,r-1).
 \end{array}
}\end{equation} The negative chirality exchanges the two spin degrees.
The labels \(A,B\) always refer to the physical unprimed/primed signs.

We measure diagram representations by weighted sums of the absolute
values of their coefficients, so that one mass controls composition as
well as the operator norm. Fix a degree \(d\) and an ordered input list
with \(|\mathbf j|=r+3-2d\), including its types. For each block
\(\nu\), let \(\mathcal D_{r,\nu}\) be the finite set of normalized bare
contraction diagrams with those inputs and that stripped output valence,
including contraction orientations. Define
\begin{equation}\protect\phantomsection\label{eq:es-diagram-mass}{\begin{gathered}
m_{r,\nu}(B)=\inf\left\{
       \sum_D|b_D|:\ 
       \widetilde B_\nu=\sum_{D\in\mathcal D_{r,\nu}}b_DD
                     \right\}, \\
m_r(B)=2m_{r,A}(B)+2m_{r,B}(B)+m_{r,C}(B).
\end{gathered}}\end{equation} Here \(\widetilde B_\nu\) is the stripped coefficient
and the infimum is taken on the span of the displayed diagrams. Absent
blocks have mass zero. All coefficients generated by
eq.~\ref{eq:es-recursion} lie in these spans: the curvature formula, the
projections, shifts and insertions use alternating contractions and
averaged slot permutations. Pairing internal alternating tensors reduces
their compositions to bare contraction diagrams of the specified final
valence.

A degree-\(d\) coefficient containing \(\ell\) scalar inputs has total
input valence
\begin{equation}\protect\phantomsection\label{eq:es-input-valence}{
 2\sum_i j_i+4(d-\ell)=2r+6-4\ell.
}\end{equation} It follows from eq.~\ref{eq:es-stripped-valences} that
the total number of bare contractions is \(1-2\ell\) in blocks \(A,B\)
and \(2-2\ell\) in block \(C\). A negative count gives the zero space.
There are consequently at most two bare contractions in any nonzero
block, independently of \(r\) and \(d\). Restoring the raw lower factors
and the isometric upper inclusions therefore gives
\begin{equation}\protect\phantomsection\label{eq:es-mass-dominates-norm}{
 \|B\|_{\mathrm{op}}\le4m_r(B).
}\end{equation} Indeed a bare diagram has norm at most two, and at most
two lower inclusions contribute an additional factor two. Taking the
infimum preserves this bound. Thus \(m_r\) is a norm on each of the
coefficient spaces under consideration.

\protect\phantomsection\label{sec:es-transfers}

\subsubsection{5.2. Linear and nonlinear
transfers}\label{linear-and-nonlinear-transfers}

Linear shifts and nonlinear insertions have different effects on the
coefficient mass. In the recursion, let \(L_iB\) denote replacement of
the \(i\)th input by its linear shift, and let \(I_i(B,C)\) denote
replacement of that input by a specified coefficient of \(Q_t\). The
estimate below holds for one fixed index \(i\); the coefficient theorem
will sum over the eligible positions.

\pandocstatementheading{\protect\phantomsection\label{lem:es-directed-wires}Lemma 5.2.}

Fix a fine assignment in the insertion of one restored bare atom into
another, followed by \(J_r^*\tau\) and extraction of a target block. The
input and output slots of every elementary operation can be placed on
distinct successive boundaries so that all identity connections are
directed from an earlier boundary to a later one. Cutting the
alternating tensors leaves only paths. Thus every closed component of
the unreduced graph contains an alternating tensor. If there are at most
\(q\) alternating tensors in both spin factors together, the product of
the closed-component factors has absolute value at most \(2^q\). The
same conclusions hold for the opened terms of \(u^cD_c\) in Lemma
\hyperref[lem:es-opening]{5.1}, with their factor \(1/p\) kept outside
the graph.

\emph{Proof.} Represent each operation on labeled tensor slots before
averaging. Its input and output are separate copies even when they have
the same valence. A fixed permutation matches these two boundaries
bijectively. Tensor products place matchings side by side, and
composition glues an output boundary only to a subsequent input
boundary. Unchanged slots are carried by identity maps.

The insertion has a definite composition order: form and restore the
inner atom, apply the restored outer atom to its selected record output,
and apply \(\tau\), \(J_r^*\) and the target extraction. The other outer
inputs and the inner derivative index are carried through the
intervening stages. Expand these stages into their elementary
operations. The following list accounts for all their slot connections.

\begin{itemize}
\tightlist
\item
  A bare atom pairs selected input slots by \(\epsilon^{ab}\) and passes
  the others to its output.
\item
  Upper restoration separates a selected slot from a symmetric tensor;
  upper-potential extraction returns the distinguished slot to the
  symmetric output. Both are averages of boundary matchings.
\item
  Multiplication and symmetrization average boundary matchings; a fine
  assignment fixes one such matching.
\item
  A raw lower restoration adjoins \(\epsilon_{ab}\), with two outgoing
  slots. A lower stripping contracts two incoming slots by
  \(\epsilon^{ab}\). All remaining slots pass forward.
\item
  The map \(\tau\) permutes slots, and \(J_r^*\) symmetrizes the new
  derivative slots into the two spin outputs as specified in Lemma
  \hyperref[lem:es-spencer]{4.1}. Each is an average of boundary
  matchings. Target upper extraction uses such a matching, and target
  lower extraction uses the stripping just described.
\end{itemize}

Number the elementary stages by \(k\). Give their input and output
boundaries heights \(2k\) and \(2k+1\), respectively, inserting identity
stages for slots carried alongside an operation. An identity connection
traverses a stage from input to output; gluing successive stages also
increases height. Cut each alternating tensor while retaining its
incident half-edges as open ends. Every remaining internal vertex has
one incoming and one outgoing identity connection, along which height
strictly increases. A closed identity component would therefore be a
directed cycle with strictly increasing height, which is impossible.

The two spin colors are preserved. Each closed component of the original
graph consequently contains an alternating tensor, and its contraction,
by the fixed alternating inverse identities, has absolute value at most
two in a two-dimensional spin factor. Distinct closed components use
disjoint alternating tensors, so there are at most \(q\) of them.
Evaluating the closed components of this unreduced graph gives the bound
\(2^q\).

For an opened term, the selected slot of the inner raw lower factor is
paired by an outer bare contraction with a slot of another outer input.
A contraction within the distinguished outer input itself is zero by
symmetry. The other outer input remains an original input of the
composite. Contracting the two incident alternating tensors gives a
signed identity from that partner slot to the new slot passed to the
symmetric output by \(u^c\). This connects an input boundary to a later
output boundary. All subsequent upper restorations and target
extractions are among the operations above. The partner normalization in
eq.~\ref{eq:es-partner-normalization} leaves the numerical factor
\(1/p\) outside this matching. Thus an opening removes two alternating
tensors and preserves the directed-boundary construction, so the same
closed-component bound applies. ∎

\pandocstatementheading{\protect\phantomsection\label{lem:es-diagram-composition}Lemma 5.3.}

Let \(D\) be a bare atom for a block \(\nu\) at order \(r-1\), with a
specified nonscalar input of order \(t\), and let \(E\) be a bare atom
for a block \(\mu\) at order \(t\) with the matching output chirality.
Suppose \(1\le t\le r-2\). Restore their one-derivative blocks and
replace the specified input of \(D\) by \(E\). Give an atom in block
\(\lambda\) the block weight \(w_\lambda\), where \(w_A=w_B=2\) and
\(w_C=1\). Then there is a representation by target bare atoms, restored
in their target blocks, such that
\begin{equation}\protect\phantomsection\label{eq:es-atomwise-insertion}{\begin{gathered}
M_rJ_r^*\tau I_i(D,E)=\sum_F\beta_FF, \\
\sum_Fw(F)|\beta_F|
 \le\frac{5C_{\mathrm w}}{t+1}
 \le\frac{K_0}{t+1}w(D)w(E), \\
C_{\mathrm w}=2^{20},\quad K_0=2^{24}.
\end{gathered}}\end{equation} The estimate is independent of the two arities, input
orders, input types and output chiralities.

\emph{Proof.} First account for the slots selected before composition.
Suppose that input \(i\) has degree \(q_i\), of which \(k_i\) slots have
already been selected, and put \(p_i=q_i-k_i\). For
\(0\le h\le\sum_i p_i\), distributing \(h\) further selections among the
residual symmetric factors gives the weights
\begin{equation}\protect\phantomsection\label{eq:es-assignment-weights}{
 \binom h{\alpha_1,\ldots,\alpha_d}
 \frac{\prod_i(p_i)_{\alpha_i}}{(\sum_i p_i)_h},
 \qquad \sum_i\alpha_i=h.
}\end{equation} The displayed weight is the total weight of the
occupancy \(\boldsymbol\alpha\). Conditional on that occupancy, average
uniformly over its \(\binom h{\alpha_1,\ldots,\alpha_d}\) allocations of
the \(h\) labeled selection slots. Each such fine assignment has weight
\(\prod_i(p_i)_{\alpha_i}/(\sum_i p_i)_h\) and specifies one slot graph.
The total weights are nonnegative and sum to one by the
falling-factorial Vandermonde identity; for \(h=0\) the weight is one. A
selection beyond the total residual valence is zero. The composite
selection on input \(i\) has exactly the normalization of a single
selection of \(k_i+\alpha_i\) slots, because
\begin{equation}\protect\phantomsection\label{eq:es-preselected-normalization}{
 (q_i)_{k_i}(q_i-k_i)_{\alpha_i}=(q_i)_{k_i+\alpha_i}.
}\end{equation} Unavailable selections have coefficient zero. Apply this
identity separately in the two spin factors. The joint assignment
weights are products of the two families of nonnegative weights, and
again sum to one. At each successive distribution, multiplying these
conditional weights preserves their total sum. A fixed fine assignment
therefore produces one slot graph with a nonnegative coefficient, and
the coefficients of all assignments sum to at most one. This remains
true when some graphs vanish. No count of the assignments, input
positions or slot permutations enters the estimate.

The operations used here have the following normalizations.
Multiplication of output polynomials is averaged symmetrization of their
free slots. Upper restoration is a normalized slot selection, and
extraction of its upper potential passes the distinguished slot to the
symmetric output. A lower stripping on a degree-\(N\) output is a
normalized single-slot selection followed by an alternating contraction,
with scalar factor \(N/(N+1)\le1\). A raw lower restoration adds the
degree-one factor \(\epsilon_{bc}u^b\). Treating that factor as a
further symmetric factor makes its derivative assignments instances of
eq.~\ref{eq:es-assignment-weights}. Finally, \(J_r^*\) averages the new
slots into the two symmetric spin outputs and \(\tau\) exchanges
specified slots. Thus every one of these operations either uses the
preceding nonnegative assignments, has scalar coefficient at most one,
or introduces an explicit alternating tensor.

For a fixed assignment and target block, all alternating tensors come
from the following operations:

\begin{center}
\begin{tabular}{@{}
  >{\raggedright\arraybackslash}p{(\linewidth - 4\tabcolsep) * \real{0.3333}}
  >{\raggedright\arraybackslash}p{(\linewidth - 4\tabcolsep) * \real{0.3333}}
  >{\raggedright\arraybackslash}p{(\linewidth - 4\tabcolsep) * \real{0.3333}}@{}}
\toprule\noalign{}
\begin{minipage}[b]{\linewidth}\raggedright
Source
\end{minipage} & \begin{minipage}[b]{\linewidth}\raggedright
Maximum number in both spin factors
\end{minipage} & \begin{minipage}[b]{\linewidth}\raggedright
Reason
\end{minipage} \\
\midrule\noalign{}
Bare contractions of \(D\) & \(2\) & eq.~\ref{eq:es-input-valence} at
order \(r-1\) \\
Bare contractions of \(E\) & \(2\) & eq.~\ref{eq:es-input-valence} at
order \(t\) \\
Raw lower restorations of \(D\) & \(2\) & At most one in each spin
factor \\
Raw lower restorations of \(E\) & \(2\) & At most one in each spin
factor \\
Stripping the target lower factors & \(2\) & At most one in each spin
factor \\
\bottomrule\noalign{}
\end{tabular}
\end{center}

The raw lower factors of the restored target atom are output data and
are excluded from its stripped graph. There are at most ten internal
alternating tensors. Lemma \hyperref[lem:es-directed-wires]{5.2} applies
to each fine assignment and bounds the closed-component contribution by
\(2^{10}\), evaluated on the unreduced graph. The same bound holds after
a lower opening, which replaces its two incident alternating tensors by
a forward boundary matching.

Reduce the open paths using the alternating inverse identities of
eq.~\ref{eq:es-spin-conventions}. Their endpoint variances give signed
identities or single alternating tensors. A path joining two slots of
one symmetric input vanishes, as does a path joining two symmetric
output slots. Every remaining path either contracts two different inputs
or passes an input slot to the output, giving a bare target atom. Its
selected-slot normalization is
eq.~\ref{eq:es-preselected-normalization}, so the reduction introduces
no additional valence factor.

For each assignment the absolute reduction coefficient is therefore at
most \(2^{10}\). The sum of the assignment weights is at most one, so
the same number bounds the total absolute coefficients in each stripped
target block. The \(++\) block is annihilated by \(M_r\), and the
remaining target weights sum to \(2+2+1=5\). Each nonzero block
multiplier of \(M_r\) has absolute value at most \(5/4\) by Lemma
\hyperref[lem:es-spencer]{4.1}. The weighted sum after \(M_r\) is
consequently at most \(2^{10}\cdot5\cdot(5/4)<2^{13}<C_{\mathrm w}\).
This proves a coefficient estimate for both spin factors together.

The preceding reduction bounds composition uniformly. To retain the
inverse-order factor, choose one lower spin factor of the inserted
coefficient and track it through the final projection. The table lists
all choices used in this argument:

\begin{center}
\begin{tabular}{@{}
  >{\raggedright\arraybackslash}p{(\linewidth - 6\tabcolsep) * \real{0.2500}}
  >{\raggedright\arraybackslash}p{(\linewidth - 6\tabcolsep) * \real{0.2500}}
  >{\raggedright\arraybackslash}p{(\linewidth - 6\tabcolsep) * \real{0.2500}}
  >{\raggedright\arraybackslash}p{(\linewidth - 6\tabcolsep) * \real{0.2500}}@{}}
\toprule\noalign{}
\begin{minipage}[b]{\linewidth}\raggedright
Inner chirality
\end{minipage} & \begin{minipage}[b]{\linewidth}\raggedright
Inner block
\end{minipage} & \begin{minipage}[b]{\linewidth}\raggedright
Chosen lower factor
\end{minipage} & \begin{minipage}[b]{\linewidth}\raggedright
Degree \(p\) of its input
\end{minipage} \\
\midrule\noalign{}
\(+\) & \(A=-+\) & unprimed & \(t+4\) \\
\(+\) & \(B=+-\) & primed & \(t\) \\
\(+\) & \(C=--\) & unprimed & \(t+4\) \\
\(-\) & \(A=-+\) & unprimed & \(t\) \\
\(-\) & \(B=+-\) & primed & \(t+4\) \\
\(-\) & \(C=--\) & unprimed & \(t\) \\
\bottomrule\noalign{}
\end{tabular}
\end{center}

There is no inner \(++\) block, and \(Q_0=0\). Hence \(t\ge1\) and
\(p\ge t\). For a double-lower block the table selects the unprimed
factor once. In the chosen factor the inserted output is
\((\iota_pf)_c\). Formula (\ref{eq:es-normalized-product}) shows that
contraction with \(u^c\) kills the unhit term. Each hit opens an
incident bare contraction of the outer atom. There are at most two such
contractions, and Lemma 5.1 gives coefficient \(1/p\) for each,
including the normalization on its partner input. Before any further
slot distribution, eq.~\ref{eq:es-partner-normalization} puts the
partner in normalized form with one fewer selected slot. The subsequent
assignment weights therefore use these updated residual valences and
again sum to one. Thus the opened coefficient sum, before the bounded
wire reductions, is at most \(2/p\) for a pair of bare atoms.

It remains to check that the final projection retains an inverse-order
factor. The outer blocks have the following signs in the chosen factor:

\begin{center}
\begin{tabular}{@{}llll@{}}
\toprule\noalign{}
Chosen factor & Outer \(A=-+\) & Outer \(B=+-\) & Outer \(C=--\) \\
\midrule\noalign{}
unprimed & lower & upper & lower \\
primed & upper & lower & lower \\
\bottomrule\noalign{}
\end{tabular}
\end{center}

These rows apply to both outer chiralities. Suppose first that the outer
sign is upper. Write its old output as \(B_a=\partial_ag/(q+1)\), where
\(q\) is its degree in \(W_{r-1}\). After inserting the chosen
coefficient into \(g\), retain its new derivative index \(c\) and call
the polynomial \(D_c\). With \(N=q+1\), the factor after \(J_r^*\tau\)
is \begin{equation}\protect\phantomsection\label{eq:es-exact-slot-map}{
 T'_a=\frac{u^c\partial_aD_c}{N},\qquad
 \deg D_c=N,\qquad N\in\{r,r+4\}.
}\end{equation} The multiplication by \(u^c\) expresses the averaged
insertion of the derivative index into the symmetric output. A nonlinear
input order satisfies \(1\le t\le r-2\), since \(Q_{r-1}\) uses only
records through order \(r-2\). In particular \(N\ge t+1\).

Set \(G=u^cD_c\). Euler's identity gives
\begin{equation}\protect\phantomsection\label{eq:es-upper-transfer-cancellation}{
 u^aT'_a=G,\qquad P_+T'=\operatorname{up}_{N+1}G.
}\end{equation} The upper coefficient therefore has the opened bound
\(2/p\le4/(t+1)\). For the lower coefficient, the mixed second
derivatives cancel:
\begin{equation}\protect\phantomsection\label{eq:es-lower-transfer-cancellation}{
 \begin{aligned}
 \partial_0(u^c\partial_1D_c)
       -\partial_1(u^c\partial_0D_c)
       &=\partial_1D_0-\partial_0D_1,\\
 \operatorname{strip}_NT'
       &=-\frac1N\operatorname{strip}_ND.
 \end{aligned}
}\end{equation} The remaining stripping is a normalized slot operation
with scalar factor \(N/(N+1)\le1\). Its diagram reduction is covered by
the preceding uniform bound, while the factor \(1/N\le1/(t+1)\) remains
explicit.

If the outer sign is lower, its factor \(\epsilon_{ba}u^b\) is
independent of the inserted input. The resulting factor is
\(\epsilon_{ba}u^bG\). Its upper projection vanishes, and its lower
coefficient is \(G\), with the same opened bound \(2/p\le4/(t+1)\). The
other spin factor undergoes the normalized operations already counted.
Together, the two tables and the upper/lower computations treat every
inner and outer block in both chiralities.

The upper and lower contributions in the chosen factor have sum at most
\(2/p+1/N\le5/(t+1)\); when the outer factor is lower, only the first
occurs. The normalized operations in the other factor and the target
weights have already been included in \(C_{\mathrm w}\). Thus the
weighted target representation satisfies
eq.~\ref{eq:es-atomwise-insertion}. Since \(5C_{\mathrm w}\le2^{24}\)
and both input block weights are at least one, the stated constant
follows. ∎

\pandocstatementheading{\protect\phantomsection\label{lem:es-transfer}Lemma 5.4.}

For \(r\ge1\), put
\begin{equation}\protect\phantomsection\label{eq:es-linear-transfer-product}{
 c_r=\frac{r(r+4)}{r(r+4)-1},\qquad
 P_r=\prod_{s=1}^rc_s,
 \qquad P_0=1.
}\end{equation} The actual recursion satisfies
\begin{equation}\protect\phantomsection\label{eq:es-transfer-bounds}{
 \begin{aligned}
 m_r(M_rJ_r^*\tau L_iB)&\le c_rm_{r-1}(B),\\
 m_r(M_rJ_r^*\tau I_i(B,C))
     &\le\frac{K_0}{t+1}m_{r-1}(B)m_t(C),
       \qquad K_0=2^{24}.
 \end{aligned}
}\end{equation} Every typed ordered quadratic coefficient of the
transferred curvature source has mass at most \(C_*(r+1)\) for a
structural constant \(C_*\). Moreover, the product converges and
\begin{equation}\protect\phantomsection\label{eq:es-transfer-product-bound}{
 1\le P_r\le P_\infty\le\exp(11/18).
}\end{equation} \emph{Proof.} Lemma
\hyperref[lem:es-diagram-composition]{5.3} gives the nonlinear estimate
on pairs of bare atoms. Choose arbitrary finite representations
\(B=\sum_Db_DD\) and \(C=\sum_Ec_EE\), with their blocks restored.
Composition is bilinear, so the atomwise representations give
\begin{equation}\protect\phantomsection\label{eq:es-representation-insertion}{
 m_r(M_rJ_r^*\tau I_i(B,C))
 \le\frac{K_0}{t+1}
       \left(\sum_Dw(D)|b_D|\right)
       \left(\sum_Ew(E)|c_E|\right).
}\end{equation} For any \(\varepsilon>0\), the two representations can
be chosen independently with weighted sums at most
\(m_{r-1}(B)+\varepsilon\) and \(m_t(C)+\varepsilon\). Letting
\(\varepsilon\) decrease to zero proves the nonlinear estimate in
eq.~\ref{eq:es-transfer-bounds}.

For a linear shift, we retain the coefficients of Lemma 5.1 rather than
the coarser reduction bound, to obtain a summable loss. Put \(n=r+4\)
and \(m=r\) for positive output chirality.

An \(A\) atom has an unprimed lower factor and a primed upper factor.
Its unchanged lower block receives the multiplier \(n/(n+1)\) from
\(M_r\); primed lower leakage has coefficient at most \(1/(m+1)\) and
receives the double-lower multiplier \(c_r\). Its column is therefore
\((n/(n+1),0,c_r/(m+1))^{\mathsf T}\).

A \(B\) atom has column \((0,m/(m+1),c_r/(n+1))^{\mathsf T}\) by the
same calculation with the spin factors exchanged. A \(C\) atom stays
double-lower with column \((0,0,c_r)^{\mathsf T}\). Hence the transfer
on the raw block masses is bounded componentwise by
\begin{equation}\protect\phantomsection\label{eq:es-linear-transfer-matrix}{
 \begin{pmatrix}
 n/(n+1)&0&0\\
 0&m/(m+1)&0\\
 c_r/(m+1)&c_r/(n+1)&c_r
 \end{pmatrix}.
}\end{equation} Its column sums with weights \((2,2,1)\) are bounded by
\(c_r\) times the input weights. The two inequalities reduce to
\begin{equation}\protect\phantomsection\label{eq:es-weighted-column-inequalities}{
 m(n+1)\le2(m+1)^2,
 \qquad n(m+1)\le2(n+1)^2.
}\end{equation} For \(n=r+4,m=r\) their right-minus-left differences are
\(r^2-r+2\) and \(r^2+15r+46\), respectively. Interchanging the spin
factors interchanges \(A,B\) and preserves their equal weights, proving
the same estimate for negative chirality.

For the quadratic curvature source, fix its ordered curvature and record
positions. First take positive output chirality, write
\(U=U_{r-1}^+(u,u')\), and put \(p=r+3\), \(p'=r-1\), \(n=r+4\),
\(m=r\). Evaluating the symmetrized direction \(b\) at \(u\otimes u'\)
describes \(J_r^*K_r\). By eq.~\ref{eq:es-curvature-decomposition}, its
\(\Psi^+=U_0^+\) contribution is, up to an overall sign,
\(\epsilon_{B'A'}u'^{B'}F_A\), where
\begin{equation}\protect\phantomsection\label{eq:es-source-weyl-polynomial}{
 F_A=\frac1{12}\epsilon^{DE}
                  (\partial_A\partial_D\Psi^+)\partial_EU.
}\end{equation} The derivatives of \(U\) sum the action on all its
unprimed slots, including derivative slots. Euler's identity and
antisymmetry give
\begin{equation}\protect\phantomsection\label{eq:es-source-normalized-atoms}{
 \begin{aligned}
 u^AF_A
 &=p\,\frac{\epsilon^{DE}(\partial_D\Psi^+)\partial_EU}{4p},\\
 \operatorname{strip}_nF
 &=\frac{p(p-1)}{n+1}\,
   \frac{\epsilon^{AF}\epsilon^{DE}
       (\partial_A\partial_D\Psi^+)\partial_F\partial_EU}
        {12p(p-1)}.
 \end{aligned}
}\end{equation} The two quotients are normalized bare atoms with one and
two contractions. In the second line, derivatives hitting \(\Psi^+\)
three times vanish under the alternating \(A,F\) contraction. The primed
factor is already a raw lower inclusion, so the two terms belong to
blocks \(B,C\). Multiplication by \(M_r\) and the weights \((2,2,1)\)
give the mass bound \(S_+(r)\) below. For the \(\Psi^-=U_0^-\)
contribution, exchange the spin factors in this calculation. Its blocks
are \(A,C\), giving \(S_-(r)\). When \(p'=0\) the latter contribution is
zero; when \(p'=1\) its double-contraction term is zero. Normalized
atoms requiring those vanishing derivatives are omitted.
\begin{equation}\protect\phantomsection\label{eq:es-source-weyl-masses}{
 \begin{aligned}
 S_+(r)&=2\frac{m}{m+1}p+c_r\frac{p(p-1)}{n+1},\\
 S_-(r)&=2\frac{n}{n+1}p'+c_r\frac{p'(p'-1)}{m+1}.
 \end{aligned}
}\end{equation} For the constant-curvature part of
eq.~\ref{eq:es-algebraic-curvature}, each spin part of
\(R_\Lambda(a,b)\) acts on the corresponding factor of \(b=u\otimes u'\)
by \(-\Lambda\eta(a,b)/6\). Insertion in the \(p+p'=2r+2\) slots
therefore gives
\begin{equation}\protect\phantomsection\label{eq:es-source-scalar-polynomial}{
 (J_r^*K_r^{\Lambda})_{AA'}
 =-\frac{\Lambda}{3}(r+1)
     \epsilon_{BA}u^B\epsilon_{B'A'}u'^{B'}U.
}\end{equation} This is a block-\(C\) atom with zero contractions and
transferred mass at most \(c_r(r+1)/3\). Since \(c_r\le5/4\), the
preceding expressions satisfy
\begin{equation}\protect\phantomsection\label{eq:es-source-uniform-bound}{
 S_+(r)\le\frac{21}{4}(r+1),\qquad
 S_-(r)\le\frac{13}{4}(r+1),\qquad
 c_r\frac{r+1}{3}\le\frac5{12}(r+1).
}\end{equation} Thus \(C_*=6\) is a valid structural choice for each
typed ordered coefficient. Negative output chirality exchanges \(n,m\)
and the two spin factors; the equal weights of \(A,B\) give the same
bound.

Finally,
\begin{equation}\protect\phantomsection\label{eq:es-summable-transfer-cost}{
 \log c_r\le\frac1{r(r+4)-1}\le\frac1{r(r+3)},\qquad
 \sum_{r\ge1}\frac1{r(r+3)}=\frac13\left(1+\frac12+\frac13\right)=\frac{11}{18}.
}\end{equation} This proves the product bound. Taking the infimum over
finite representations of \(B\) in the linear column estimates gives the
first line of eq.~\ref{eq:es-transfer-bounds}. ∎

Define the rescaled mass
\begin{equation}\protect\phantomsection\label{eq:es-rescaled-mass}{
 \gamma_r=4\frac{P_\infty}{P_r}m_r.
}\end{equation} This rescaling absorbs the accumulated linear loss: it
dominates the operator norm and makes each specified linear transfer
nonexpanding. For a nonlinear insertion, the factor relative to the two
input masses is
\begin{equation}\protect\phantomsection\label{eq:es-insertion-rescaling}{
 \frac{P_{r-1}P_t}{4P_\infty P_r}\le\frac14.
}\end{equation} The product rule sums the two possible intermediate
record chiralities. A scalar \(\Lambda\) position has derivative zero.
Absorbing that sum gives one structural constant \(K\le K_0\) for the
rescaled insertion bound \(K/(t+1)\). The arity induction below counts
the eligible ordered positions.

\Needspace{10\baselineskip}
\protect\phantomsection\label{sec:es-ordered-coefficients}

\subsubsection{5.3. The coefficient
theorem}\label{the-coefficient-theorem}
\leavevmode\par\nobreak

\pandocstatementheading{\protect\phantomsection\label{thm:es-coefficients}Theorem 5.5.}

There are structural constants \(a,b>0\) such that every typed ordered
coefficient in eq.~\ref{eq:es-ordered-expansion} satisfies
\begin{equation}\protect\phantomsection\label{eq:es-coefficient-budget}{
 \gamma_r(B^{[d]}_{r,\mathbf j}),\ 
 \|B^{[d]}_{r,\mathbf j}\|_{\mathrm{op}}
 \le ab^{d-2}r\frac{r!}{\prod_i j_i!},
 \qquad |\mathbf j|=r+3-2d,\quad d\ge2.
}\end{equation} \emph{Proof.} For degree two, a curvature source
introduced at order \(s\) has ordered indices \((0,s-1)\) and rescaled
mass at most \(C(s+1)\). Subsequent linear shifts distribute their
increments among the two existing positions. At final indices
\(j+k=r-1\), their total contribution is bounded by
\begin{equation}\protect\phantomsection\label{eq:es-quadratic-sum}{
 C\sum_{s=1}^{r-j}(s+1)\binom{r-s}{j}
 =C\left\{\binom{r+1}{j+2}+\binom r{j+1}\right\}.
}\end{equation} The identity follows by summing binomial coefficients,
or by writing \(s+1=2+\sum_{\ell=1}^{s-1}1\) and applying the
hockey-stick identity twice. Multiplication by \(j!k!/r!\) bounds this
by
\begin{equation}\protect\phantomsection\label{eq:es-quadratic-normalization}{
 C\left\{\frac{r+1}{(j+1)(j+2)}+\frac1{j+1}\right\}\le2Cr.
}\end{equation} The reverse source ordering has the same bound with
\(j,k\) exchanged. Choose \(a\) to include its finite multiplicity and
the fixed source types.

Normalize each rescaled coefficient mass by multiplying by
\(\prod_i j_i!/r!\). Induct on the arity and, within each arity, on
\(r\). Suppose a normalized degree-\(d\) coefficient is bounded by
\(A_dr\). Set \(k=|\mathbf j|=r+3-2d\). A linear shift into position
\(i\) comes from index \(j_i-1\) and contributes at most
\(A_d(r-1)j_i/r\). Summing over the positions gives at most \(A_dk\).

For a nonlinear insertion, let the inner and outer arities be
\(e,f\ge2\), so that \(e+f=d+1\). The order of the old outer position is
\begin{equation}\protect\phantomsection\label{eq:es-inserted-order}{
 t=\sum_{\mathrm{inner}}j_i+2e-3\ge1.
}\end{equation} The outer factorial denominator \(t!\) cancels the inner
factorial numerator. After normalization by \(\prod_i j_i!/r!\), Lemma
\hyperref[lem:es-transfer]{5.4} and eq.~\ref{eq:es-insertion-rescaling}
give the bound
\begin{equation}\protect\phantomsection\label{eq:es-normalized-insertion}{
 K A_fA_e\frac{(r-1)t}{r(t+1)}\le K A_fA_e.
}\end{equation} For a fixed final ordered list and fixed \(e,f\), there
are \(f\) possible outer positions. At each position, the inner inputs
are determined as its consecutive ordered sublist. Thus the product rule
contributes \(f\), with no further partition factor. The scalar
positions contribute zero to differentiation. Since \(r-k=2d-3\), the
target bound leaves \((2d-3)A_d\) for the nonlinear terms. We therefore
close the induction with
\begin{equation}\protect\phantomsection\label{eq:es-arity-recursion}{
 A_2=a,\qquad
 A_d=\frac K{2d-3}
       \sum_{\substack{e+f=d+1\\e,f\ge2}}fA_eA_f
 \le\frac{2K}{3}
       \sum_{\substack{e+f=d+1\\e,f\ge2}}A_eA_f.
}\end{equation} For the last inequality, interchange \(e,f\) in the sum
and use \((d+1)/(2(2d-3))\le2/3\) for \(d\ge3\). The Catalan recurrence
then yields
\begin{equation}\protect\phantomsection\label{eq:es-catalan-bound}{
 A_d\le a(2Ka/3)^{d-2}\operatorname{Cat}_{d-2}
       \le ab^{d-2},\qquad b=8Ka/3.
}\end{equation} The rescaled mass dominates the operator norm by
eq.~\ref{eq:es-mass-dominates-norm} and eq.~\ref{eq:es-rescaled-mass}. ∎

Group the input types using the fixed direct-sum norms in
eq.~\ref{eq:es-record-norm}. There are at most \(3^d\) ordered
input-type lists and two output chiralities. Hence the grouped
coefficients satisfy the same estimate with
\begin{equation}\protect\phantomsection\label{eq:es-grouped-constants}{
 a_0=18a,\qquad b_0=3b,
 \qquad 2\cdot3^dab^{d-2}=a_0b_0^{d-2}.
}\end{equation}

\pandocstatementheading{\protect\phantomsection\label{cor:es-nonlinear-field}Corollary 5.6.}

Write \(Q_a(X)\) for the nonlinear record vector with scalar component
zero, and set
\begin{equation}\protect\phantomsection\label{eq:es-nonlinear-field-norm}{
 \|Q(X)\|_s=\sup_{\|a\|\le1}\|Q_a(X)\|_{\mathcal B_s}.
}\end{equation} Let \(M=\|X\|_R\), let \(0<s<R\), and suppose
\(\vartheta=b_0s^2M<1\). Then the nonlinear part of the original record
derivation converges in \(\mathcal B_s\) and satisfies
\begin{equation}\protect\phantomsection\label{eq:es-nonlinear-bound}{
 \|Q(X)\|_s
 \le a_0sM^2\left(2+\frac{s}{R-s}\right)
                    \frac{2-\vartheta}{(1-\vartheta)^2}.
}\end{equation} It is a holomorphic map on every open input ball for
which these inequalities hold uniformly with strict margins.

\emph{Proof.} Put \(c_j=R^j\|X_j\|/j!\) and \(q=s/R\). In the
degree-\(d\) sum, \(r=\sum_i j_i+2d-3\). The unweighted product sum is
at most \(M^d\), while
\begin{equation}\protect\phantomsection\label{eq:es-weighted-product-bound}{
 \sum_{j_1,\ldots,j_d}\left(\sum_i j_i\right)
        q^{\sum_i j_i}\prod_i c_{j_i}
 \le dM^d\frac{q}{1-q}.
}\end{equation} Indeed \(j q^j\le q/(1-q)\), and one sums the
distinguished position first. Applying
eq.~\ref{eq:es-coefficient-budget} with the grouped constants bounds the
degree-\(d\) contribution by \(a_0sM^2d\vartheta^{d-2}(2+q/(1-q))\).
Summing the identity
\(\sum_{d\ge2}d\vartheta^{d-2}=(2-\vartheta)/(1-\vartheta)^2\) gives
eq.~\ref{eq:es-nonlinear-bound}. Finite input and output truncations are
continuous polynomial maps. The same estimates, with a strict
output-radius margin, make them converge locally uniformly in the output
norm. Their limit is therefore holomorphic. ∎

\protect\phantomsection\label{sec:es-flows}

\subsection{6. Analytic record flows}\label{analytic-record-flows}

We integrate the horizontal derivations on decreasing record radii. A
scalar majorant controls the radius loss and yields a flow holomorphic
in the initial record and direction. The canonical inclusions
\(\iota_{u,v}:\mathcal B_u\to\mathcal B_v\), \(0<v<u\), are those of §2
and have norm at most one. The horizontal vector field is
\begin{equation}\protect\phantomsection\label{eq:es-horizontal-vector-field}{
 \mathsf H_a(X)_r=i_aU_{r+1}+Q_{r,a}(X),\qquad
 \mathsf H_a\Lambda=0.
}\end{equation} The leading shift has norm at most \(\|a\|/(u-v)\) from
\(\mathcal B_u\) to \(\mathcal B_v\). Corollary
\hyperref[cor:es-nonlinear-field]{5.6} makes the full field holomorphic
between suitable smaller-radius spaces on bounded input balls. The three
radii in the theorem below distinguish the input space, the space
containing the flow, and the smaller space where its equation and
compositions are compared. The spatial normal-coordinate neighborhood is
constructed in §7.

\pandocstatementheading{\protect\phantomsection\label{thm:es-flow}Theorem 6.1.}

For every \(R,M,A>0\), there are \(0<\sigma<s<R\) and \(T,C>0\) with the
following properties. On an open neighborhood
\(\mathcal U\subset V\times\mathcal B_R\) of \[
 \{(a,X):\|a\|\le A,\ \|X\|_R\le M\}
\] there is a jointly holomorphic map
\begin{equation}\protect\phantomsection\label{eq:es-flow-map}{\begin{gathered}
\Phi:\{t\in\mathbb C:|t|<T\}\times\mathcal U
                 \longrightarrow\mathcal B_s, \\
\Phi(t,a,X)=\Phi_a^t(X), \\
\|\Phi_a^t(X)\|_s\le C.
\end{gathered}}\end{equation} It is initialized by \(\Phi_a^0(X)=\iota_{R,s}X\) and
satisfies
\begin{equation}\protect\phantomsection\label{eq:es-flow-equation}{
 \iota_{s,\sigma}\,\partial_t\Phi_a^t(X)
       =\mathsf H_a\bigl(\Phi_a^t(X)\bigr)
       \quad\text{in }\mathcal B_\sigma.
}\end{equation} The right side is the original polynomial record
derivation, evaluated by its convergent sequence expansion. Its
coordinate equations determine the analytic-time solution uniquely.

The initial-value derivative \(Y_h(t)=D_X\Phi_a^t(X)[h]\) is a
holomorphic \(\mathcal B_s\)-valued function for \(h\in\mathcal B_R\).
On the displayed bounded input set it has a common bound
\(\|Y_h(t)\|_s\le C\|h\|_R\), after increasing \(C\), and
\begin{equation}\protect\phantomsection\label{eq:es-flow-variational-equation}{
 \begin{aligned}
 Y_h(0)&=\iota_{R,s}h,\\
 \iota_{s,\sigma}\,\partial_tY_h(t)
   &=D\mathsf H_a\bigl(\Phi_a^t(X)\bigr)[Y_h(t)]
                         \quad\text{in }\mathcal B_\sigma.
 \end{aligned}
}\end{equation} After choosing \(T\) smaller, compositions initialized
on the intermediate space \(\mathcal B_s\) are defined and satisfy, in
the common space \(\mathcal B_\sigma\),
\begin{equation}\protect\phantomsection\label{eq:es-flow-composition}{
 \begin{aligned}
 \Phi_a^t\bigl(\Phi_a^u(X)\bigr)
       &=\iota_{s,\sigma}\Phi_a^{t+u}(X),\\
 \Phi_a^{-t}\bigl(\Phi_a^t(X)\bigr)&=\iota_{R,\sigma}X,
 \end{aligned}
}\end{equation} whenever \(|t|,|u|,|t+u|<T\). The first expression uses
the flow with initial scale \(s\), followed by its inclusion in
\(\mathcal B_\sigma\).

Every specified finite word of horizontal flows and bounded full-spin
actions has a common positive record radius, time domain and bound, also
for its inverse word. These maps and their initial-value derivatives are
jointly holomorphic. The coefficient involution of Definition
\hyperref[def:es-lorentz-involution]{2.4} satisfies
\begin{equation}\protect\phantomsection\label{eq:es-flow-reality}{
 \mathfrak c\Phi_a^t(X)
   =\Phi_{\mathfrak ca}^{\bar t}(\mathfrak cX).
}\end{equation} Consequently Lorentz-real directions and real times,
together with Lorentz-real full-spin actions, preserve the fixed locus.

\emph{Proof.} The scalar majorant will give convergence and joint
holomorphy; coefficientwise uniqueness will then give the composition
laws. Choose larger bounds \(\widehat M=M+1\) and \(\widehat A=A+1\). We
construct the flow on the product of the open balls with these bounds,
retaining an open neighborhood of the original bounded set under all
later restrictions.

Use the recursively ordered expansion of §5, grouped in the fixed
direct-sum norms:
\begin{equation}\protect\phantomsection\label{eq:es-grouped-expansion}{
 Q_{r,a}(X)=\sum_{d\ge2}\
 \sum_{\substack{\mathbf j\in\mathbb N_0^d\\
                  |\mathbf j|=r+3-2d}}
 B^{[d]}_{r,\mathbf j}(a;X_{j_1},\ldots,X_{j_d}).
}\end{equation} Theorem \hyperref[thm:es-coefficients]{5.5} and
eq.~\ref{eq:es-grouped-constants} give structural constants
\(a_0,b_0>0\) such that
\begin{equation}\protect\phantomsection\label{eq:es-grouped-budget}{
 \|B^{[d]}_{r,\mathbf j}(a;\cdot)\|
 \le\|a\|a_0b_0^{d-2}r\frac{r!}{\prod_i j_i!}.
}\end{equation} These constants include the summable linear transfer
loss:
\begin{equation}\protect\phantomsection\label{eq:es-transfer-product}{
 c_r=\frac{r(r+4)}{r(r+4)-1},\qquad
 \prod_{r\ge1}c_r\le e^{11/18}.
}\end{equation} For formal time series with initial norm at most
\(\widehat M\) and \(\|a\|\le\widehat A\), divide the record components
by \(r!\). They are majorized by the nonnegative scalar recursion
\begin{equation}\protect\phantomsection\label{eq:es-flow-majorant}{\begin{gathered}
\partial_\tau F
   =\widehat A\left(F_z+z\frac{\mathrm d}{\mathrm dz}V(z,F)\right), \\
V(z,w)=\frac{a_0zw^2}{1-b_0z^2w}, \\
F(0,z)=\frac{\widehat M}{1-z/R}.
\end{gathered}}\end{equation} Here \(z\) is the generating variable for derivative
order and \(\tau\) is the scalar majorant time. The derivative of
\(V(z,F)\) is total. Indeed write \(X_r(t)=\sum_{k\ge0}X_{r,k}t^k\) and
\(F(\tau,z)=\sum_{k,r\ge0}f_{k,r}\tau^kz^r\). Initially
\(\|X_{r,0}\|/r!\le\widehat M R^{-r}=f_{0,r}\). The leading shift is
bounded by the corresponding coefficient of \(\widehat A F_z\). If the
record coefficients through time order \(k\) are dominated, the
degree-\(d\) nonlinear term is bounded by
\begin{equation}\protect\phantomsection\label{eq:es-flow-majorization}{
 \begin{aligned}
 &\widehat A a_0b_0^{d-2}r
 [\tau^k]\!
 \sum_{j_1+\cdots+j_d=r+3-2d}
       \prod_{i=1}^d[z^{j_i}]F(\tau,z)\\
 &\qquad= [\tau^kz^r]\,\widehat A z\frac{\mathrm d}{\mathrm dz}
       \left(a_0b_0^{d-2}z^{2d-3}F^d\right).
 \end{aligned}
}\end{equation} For fixed \(r\) the sum is finite. The coefficient
equations for \(X_{r,k+1}\) and \(f_{k+1,r}\) then give
\(\|X_{r,k+1}\|/r!\le f_{k+1,r}\). Induction in \(k\), simultaneously
for every \(r\), proves formal coefficientwise domination. Each
\(X_{r,k}\) is a polynomial in \(a\) and finitely many initial record
components.

Put \(W_*=4\widehat M+1\) and choose \(0<\rho\le R/8\) such that
\(16b_0\rho^2W_*\le1\). The finite characteristic equations are
\begin{equation}\protect\phantomsection\label{eq:es-characteristics}{
 \dot z=-\widehat A(1+zV_w),\qquad
 \dot w=\widehat A zV_z,\qquad
 z(0)=\zeta,\quad w(0)=\widehat M/(1-\zeta/R).
}\end{equation} On \(|z|\le2\rho\), \(|w|\le W_*\) their denominator is
bounded away from zero, and \[
 |zV_w|\le16a_0\rho^2W_*,\qquad
 |zV_z|\le5a_0\rho W_*^2.
\] The vector field and its first derivatives are bounded in the scaled
coordinates \((z/\rho,w/W_*)\). Its finite integral equation is a
contraction for a common sufficiently small complex time disc, with
scaled displacement less than \(1/8\). For \(|\zeta|\le\rho\) the
initial second scaled coordinate is less than \(2/7\), so these
solutions remain in the displayed parent domain. On \(|\zeta|=3\rho/4\)
the estimate \(|z(\tau,\zeta)-\zeta|\le\rho/8\) and Rouché's theorem
give a unique simple inverse for \(|z_*|<\rho/2\). Hence
\(F(\tau,z_*)=w(\tau,\zeta(\tau,z_*))\) is jointly holomorphic and
bounded by \(W_*\) there. Its Taylor series solves the nonnegative
formal recursion above.

Take \(s=\rho/4\) and restrict to a smaller positive time disc.
Coefficientwise domination now gives normal convergence of the double
record series in \(\mathcal B_s\), with a uniform bound \(C_0=2W_*\).
Its finite time and output truncations are continuous polynomial maps of
\((a,X)\) using finitely many initial records. On strictly smaller time
discs they converge in this norm uniformly on the enlarged input balls.
The scalar double-series tails tend uniformly to zero. Thus the limit is
jointly Fréchet-holomorphic in \((t,a,X)\) with values in
\(\mathcal B_s\).

Choose \(0<\sigma<s/2\) with \(b_0\sigma^2(C_0+1)<1/2\). Corollary
\hyperref[cor:es-nonlinear-field]{5.6}, applied from \(\mathcal B_s\) to
\(\mathcal B_\sigma\), and the shift bound make \(\mathsf H_a\) a
holomorphic map on the ball \(\|Z\|_s<C_0+1\) into
\(\mathcal B_\sigma\). The normally convergent time series can be
differentiated in \(\mathcal B_s\) on smaller time discs. Each
coordinate of its derivative equals the corresponding finite polynomial
coefficient of \(\mathsf H_a(\Phi_a^t(X))\). Both sequences belong to
\(\mathcal B_\sigma\), which proves eq.~\ref{eq:es-flow-equation} with
the original vector field.

For uniqueness, each coordinate of an analytic-time solution satisfies
the finite polynomial equation for that coordinate. Differentiating
successively at zero determines its time coefficients by
\begin{equation}\protect\phantomsection\label{eq:es-flow-time-recursion}{
 (k+1)X_{r,k+1}=[t^k](\mathsf H_aX(t))_r.
}\end{equation} For fixed \(r,k\), the right side uses finitely many
previously determined time coefficients and record components. Two
initialized analytic solutions therefore have identical coordinate
Taylor series and agree near zero, hence throughout a connected common
time domain. Solutions obtained from different input radii agree after
their canonical inclusions whenever both are defined.

Joint holomorphy and the enlarged input ball give the initial-value
Fréchet derivative into \(\mathcal B_s\). Cauchy's formula on a
fixed-radius initial-data disc supplies its uniform operator bound on
\(\|X\|_R\le M\), \(\|a\|\le A\). The chain rule applied to
eq.~\ref{eq:es-flow-equation} in \(\mathcal B_\sigma\) proves
eq.~\ref{eq:es-flow-variational-equation}. Here the differential is
taken for the holomorphic map
\(\mathsf H_a:\mathcal B_s\to\mathcal B_\sigma\). Increasing \(C\)
covers both value and derivative bounds.

For composition, apply the construction a second time with input scale
\(s\), initial bound \(C_0+1\) and direction bound \(\widehat A\). It
gives a positive output radius \(s_1<s\) and a common time disc for all
intermediate records \(\Phi_a^u(X)\). Decrease \(\sigma\) below \(s_1\)
and keep the preceding vector-field margin. Take \(T\) smaller than both
construction times. The curves \(t\mapsto\Phi_a^{t+u}(X)\) and
\(t\mapsto\Phi_a^t(\Phi_a^u(X))\) solve the same initialized coordinate
equations. Their common time domain \(|t|<T\), \(|t+u|<T\) is connected
and contains zero. The uniqueness just proved gives their equality in
\(\mathcal B_\sigma\). Taking \(u=t\) and the outer time \(-t\) gives
the inverse identity with \(\iota_{R,\sigma}X\). This proves
eq.~\ref{eq:es-flow-composition}.

For a bounded set of full-spin matrices, their inverses are also
bounded, since each spin matrix is two-dimensional with determinant one.
Choose a common spin-factor norm bound \(L\ge1\) for these matrices and
their inverses. The action on \(W_r^\pm\) has norm at most \(L^{2r+4}\),
so it maps a bounded \(\mathcal B_u\) ball to a bounded \(\mathcal B_v\)
ball for \(0<v<u/L^2\). Finite truncations and this strict margin prove
joint holomorphy in the spin element and the record. Apply each
horizontal or spin operation to the output radius and bound of the
preceding one. A specified finite word and its inverse require finitely
many such restrictions, leaving a positive common radius and time
domain. Differentiating these holomorphic maps gives the corresponding
finite-word derivatives.

Finally the polynomial tensor identities imply
\(\mathfrak c(\mathsf H_aX)=\mathsf H_{\mathfrak ca}(\mathfrak cX)\).
Conjugating the convergent time series and applying initialized
uniqueness gives eq.~\ref{eq:es-flow-reality}. The fixed-locus
conclusion follows. ∎

\protect\phantomsection\label{sec:es-realization}

\subsection{7. Einstein realization and normal
uniqueness}\label{einstein-realization-and-normal-uniqueness}

The radial Jacobi equation provides the passage from curvature records
to a metric. Its formal solution determines the normal metric and its
Einstein identities from the completed curvature list. The horizontal
flow makes the radial curvature series converge; solving the same Jacobi
equation analytically then gives a metric with that formal Taylor
series. Write \(D_\delta=\{x\in V:|x^j|<\delta,\ 1\le j\le4\}\) in the
fixed marked tangent basis.

\pandocstatementheading{\protect\phantomsection\label{prop:es-normal-uniqueness}Proposition 7.1.}

The completed ordered curvature list determines a unique formal metric
in marked geodesic normal coordinates. Its curvature derivatives are the
prescribed completion and its Ricci tensor is \(\Lambda g\). The
construction commutes with extension of the complex coefficient algebra.
Two marked analytic Einstein metrics with equal free records and
numerical constant are equal in these normal coordinates.

\emph{Proof.} With \(x\in V\) a formal spatial variable, set
\begin{equation}\protect\phantomsection\label{eq:es-formal-radial-curvature}{
 R_x(t)=\sum_{r\ge0}\frac{t^r}{r!}
                  \mathcal C_r(x,\ldots,x;\,\cdot).
}\end{equation} Solve the formal Jacobi equation
\begin{equation}\protect\phantomsection\label{eq:es-formal-jacobi}{\begin{gathered}
J_x''(t;v)+R_x(t)(J_x(t;v),x)x=0, \\
J_x(0;v)=0,\quad J_x'(0;v)=v, \\
g_x(v,w)=\eta(J_x(1;v),J_x(1;w)).
\end{gathered}}\end{equation} The coefficient of \(J_x\) has spatial order at least
two, so successive spatial degrees determine \(J_x\) uniquely. At each
degree the dependence on \(t\) is polynomial, making evaluation at
\(t=1\) formally well-defined. At each fixed spatial degree, the
recursion uses finitely many completed derivatives and divides only by
positive integers. Each metric coefficient is therefore polynomial in a
finite record prefix and \(\Lambda\), as are the inverse-metric
coefficients obtained by formal inversion at \(g_0=\eta\). The Einstein
residual \(\operatorname{Ric}(g)-\Lambda g\), the normal-coordinate
residual \(\Gamma^k_{ij}(x)x^ix^j\), and each discrepancy between a
curvature derivative at the origin and its prescribed completion all
have coefficients of this kind.

For any one such coefficient, Lemma
\hyperref[lem:es-finite-witness]{A.1} gives an analytic Einstein metric
realizing a finite Lorentz-real prefix long enough to contain every
record used by it. In the parallel frame along its radial geodesic
\(\gamma_x\), its curvature satisfies \[
 \left.\frac{\mathrm d^r}{\mathrm dt^r}
       P_{t\to0}R_{\gamma_x(t)}\right|_{t=0}
     =(\nabla_x^rR)_0=\mathcal C_r(x,\ldots,x;\,\cdot).
\] The witness metric in exponential coordinates is recovered by its
Jacobi fields through eq.~\ref{eq:es-formal-jacobi}. Its Taylor
coefficients through the required degree are consequently those just
constructed, and the residual coefficient is zero. Polynomial
continuation from the maximal real form, as in Proposition
\hyperref[prop:es-completion]{3.2}, proves the identity on independent
complex records. This gives all prescribed curvature derivatives, the
Einstein equation and marked normal coordinates. These polynomial
identities remain valid over every commutative complex coefficient
algebra.

Conversely, any formal metric with this completed curvature list,
expressed in marked normal coordinates and a radial parallel frame, has
radial curvature eq.~\ref{eq:es-formal-radial-curvature}. Its Jacobi
equation recovers its metric by eq.~\ref{eq:es-formal-jacobi}, proving
uniqueness. Equal Taylor series give equal analytic germs. An isometry
with the prescribed point and tangent map intertwines their exponential
maps, so it is uniquely determined as well. ∎

\pandocstatementheading{\protect\phantomsection\label{thm:es-realization}Theorem 7.2.}

Every analytic record \(X\) determines a marked holomorphic Einstein
four-germ by the convergent horizontal flow and a radial Jacobi
equation. Its entire free record is the original one. For bounded
holomorphic input families, its normal coordinates and radial spin
frames have common positive domains. The construction is unique up to
the prescribed marked isometry. Conversely, the free records of every
marked holomorphic Einstein four-germ satisfy a factorial growth bound.

\emph{Proof.} Fix \(R,M>0\) with \(\|X\|_R\le M\). Apply Theorem
\hyperref[thm:es-flow]{6.1} on the larger open input ball of radius
\(M+1\), with direction bound one. It gives a common time \(T>0\),
output radius \(s>0\) and output bound. Fix \(0<2\varepsilon<T\). For
\(\|a\|<\varepsilon\), put
\begin{equation}\protect\phantomsection\label{eq:es-rescaled-flow}{
 \operatorname{Fl}_{\mathsf H_a}^{\,t}X
       :=\Phi_{a/\varepsilon}^{\varepsilon t}(X),
       \qquad |t|<2.
}\end{equation} Linearity of \(\mathsf H_a\) in \(a\) and initialized
uniqueness in Theorem \hyperref[thm:es-flow]{6.1} identify this with the
original short-time flow wherever both are defined, and make it
independent of \(\varepsilon\) on common smaller record spaces. It is
jointly holomorphic in \((t,a,X)\) and uniformly bounded there. Define
the algebraic-curvature-valued holomorphic map
\begin{equation}\protect\phantomsection\label{eq:es-analytic-radial-curvature}{
 \widehat R_X(y)=\mathcal C_0\bigl(
           \operatorname{Fl}_{\mathsf H_y}^{\,1}X\bigr),
           \qquad \|y\|<\varepsilon.
}\end{equation} For \(\|x\|<\varepsilon/2\) and \(|t|<2\), initialized
uniqueness gives \(\operatorname{Fl}_{\mathsf H_{tx}}^{\,1}X
 =\operatorname{Fl}_{\mathsf H_x}^{\,t}X\). The last identity of
eq.~\ref{eq:es-cartan-identities}, iterated on the polynomial completed
tensors, identifies all time derivatives at zero. Hence
\begin{equation}\protect\phantomsection\label{eq:es-convergent-radial-curvature}{
 \widehat R_X(tx)
   =\mathcal C_0\bigl(\operatorname{Fl}_{\mathsf H_x}^{\,t}X\bigr)
   =\sum_{r\ge0}\frac{t^r}{r!}
                   \mathcal C_r(x,\ldots,x;\,\cdot).
}\end{equation} This series converges on the displayed time disc and
realizes the formal radial series of Proposition
\hyperref[prop:es-normal-uniqueness]{7.1}.

Solve the analytic Jacobi equation by its integral form
\begin{equation}\protect\phantomsection\label{eq:es-analytic-jacobi}{\begin{gathered}
J_x(t;v)=tv-
    \int_0^t(t-u)\widehat R_X(ux)(J_x(u;v),x)x\,\mathrm du, \\
g_X(x)(v,w)=\eta(J_x(1;v),J_x(1;w)).
\end{gathered}}\end{equation} The flow bound gives a uniform bound for
\(\widehat R_X\). The norm of the linear integral operator on the right,
on \(|t|\le3/2\), is at most \(C\|x\|^2\). A common smaller
\(x\)-neighborhood makes this norm less than \(1/2\). The Picard series
therefore converges uniformly, is holomorphic in \((t,x,X)\), and gives
\(J_x(1;\,\cdot)=I+O(\|x\|^2)\) uniformly on the record ball. Thus
\(g_X\) is a holomorphic nondegenerate metric with \(g_X(0)=\eta\) and
\(\partial g_X(0)=0\).

Its spatial Taylor series solves the formal Jacobi recursion in
Proposition \hyperref[prop:es-normal-uniqueness]{7.1}. By uniqueness of
that recursion, it is the formal normal Einstein metric of that
proposition. All Taylor coefficients of the analytic Einstein and
normal-coordinate residuals consequently vanish, and every ordered
curvature derivative at the mark is \(\mathcal C_r\). The identity
theorem gives \(\operatorname{Ric}(g_X)=\Lambda g_X\) and normal
coordinates, while the free projection gives the prescribed records.
Proposition \hyperref[prop:es-normal-uniqueness]{7.1} proves marked
uniqueness. The common domain and holomorphic dependence just obtained,
followed by radial spin parallel transport, give the family assertion;
Proposition \hyperref[prop:es-uniform-realization]{7.3} records the
uniform bounds.

For the converse, take an analytic Einstein metric and
parallel-transport its Weyl spinors along geodesics with initial
velocities \(\xi\otimes\pi\), where \(\xi,\pi\) are unit spinors.
Compactness of these unit spinor pairs gives a common affine time disc
and curvature bound for the finite holomorphic geodesic and frame
equations. Cauchy's formula bounds the coherent evaluations of the free
order-\(r\) record by \(Cr!R_0^{-r}\). Lemma
\hyperref[lem:es-record-growth]{2.3} converts this to the required
tensor-norm bound after a radius decrease. ∎

\pandocstatementheading{\protect\phantomsection\label{prop:es-uniform-realization}Proposition 7.3.}

For every \(R,M>0\) there are \(\delta>0\) and \(C<\infty\) such that
the normal metric of Theorem \hyperref[thm:es-realization]{7.2} is
defined on \(D_{2\delta}\subset V\) for every \(\|X\|_R\le M\) and
satisfies
\begin{equation}\protect\phantomsection\label{eq:es-uniform-metric-bound}{
 \sup_{D_{2\delta}}(\|g_X\|+\|g_X^{-1}\|)\le C.
}\end{equation} On an open neighborhood of this record ball the maps
\(X\mapsto g_X|_{D_\delta}\) and \(X\mapsto g_X^{-1}|_{D_\delta}\) are
Fréchet-holomorphic maps into the Banach spaces of bounded holomorphic
tensor-valued functions. The normal-coordinate maps, their inverses and
radial spin frames have the same property after common restrictions.
Every fixed finite spatial derivative has a uniform interior bound.

\emph{Proof.} The enlarged input ball and common flow bounds in Theorem
\hyperref[thm:es-realization]{7.2} allow us to choose \(\delta\) so that
the Jacobi integral operator is a contraction on a neighborhood of
\(\overline{D_{4\delta}}\). The Picard series then converges uniformly
on the product domain and is holomorphic in the Banach parameter \(X\).
For the endpoint operator \(L_X(x)v=J_x(1;v)\), decrease \(\delta\)
until \(\|L_X-I\|<1/2\) uniformly. The normally convergent Neumann
series bounds \(L_X^{-1}\) there. Since \[
 g_X=L_X^{\mathsf T}\eta L_X,\qquad
 g_X^{-1}=L_X^{-1}\eta^{-1}(L_X^{-1})^{\mathsf T},
\] both tensors have common bounds and holomorphic parameter dependence
on the larger spatial parent.

Cauchy's formula across the remaining spatial buffer bounds every
specified derivative on \(D_{2\delta}\). The Christoffel and
spin-connection coefficients are therefore uniformly bounded on smaller
parents. The finite geodesic and parallel-spin integral equations have
uniformly convergent Picard series there, giving common
normal-coordinate and radial-frame maps with holomorphic parameter
dependence. Each based exponential map has derivative the identity in
its specified tangent frame at zero. Restricting its endpoint and
initial-velocity domains makes the derivative uniformly close to that
identity; the contraction proof of the finite inverse theorem gives a
common inverse and its holomorphic dependence.

Finally expand each of these bounded maps in the spatial variables on a
larger polydisc. Its coefficients are holomorphic in \(X\), and the
uniform spatial bound makes its Taylor polynomials converge locally
uniformly in the bounded-function norm on a smaller polydisc. Their
limit is Fréchet-holomorphic in that norm. Apply this with a final
common restriction to \(D_\delta\). ∎

\protect\phantomsection\label{sec:es-naturality}

\subsection{8. Point and spin-frame changes and Lorentz
reality}\label{point-and-spin-frame-changes-and-lorentz-reality}

Point and frame changes give different marked realizations of the same
local geometry. Let
\(G=\operatorname{Spin}(V,\eta)=\operatorname{SL}(S)\times\operatorname{SL}(S')\).
We write \(A\) also for the induced tangent map of \(A\in G\), and
\(A^*U\) for insertion of its two spin maps in every argument of the
records. A full spin lift is retained as part of a comparison. The
normal realization of a record \(U\) will be denoted by \(g[U]\). The
coefficient \(\Lambda\) is suppressed in this notation and is unchanged
by each comparison. The marked exponential map of this realization
identifies its normal neighborhood with a neighborhood of \(0\in V\).

Use the rescaled flow of eq.~\ref{eq:es-rescaled-flow}, on a common
bounded record ball. Take \(\|a\|<\varepsilon\) small enough that \(ta\)
lies in the common normal parent for \(|t|<2\).

Put
\begin{equation}\protect\phantomsection\label{eq:es-point-frame-action}{
 x_a=\operatorname{Exp}^{g[U]}_o(a),\qquad
 \nu_a:V\longrightarrow T_{x_a}M,
 \qquad
 \mathcal R_{a,A}U
 =A^*\operatorname{Fl}_{\mathsf H_a}^{\,1}U,
}\end{equation} where \(\nu_a\) is the specified spin parallel transport
along the radial geodesic from \(o\) to \(x_a\).

A short record comparison from \(U\) to \(U'\) is a pair \((a,A)\) with
\(U'=\mathcal R_{a,A}U\), together with representatives on which the
displayed flow and normal maps are defined. A finite diagram of such
comparisons is conditioned if its spin maps and inverse spin maps are
bounded, and all exponential maps, inverse maps and specified finite
composites have images in their assigned normal parents with uniformly
positive buffer and inverse margins on the bounded family. These
conditions are imposed on representatives; taking a common restriction
preserves the represented local comparisons.

\pandocstatementheading{\protect\phantomsection\label{thm:es-physical-arrows}Theorem 8.1.}

Short record comparisons are in bijection with short local isometries
between the normal realizations with their specified full spin lifts.
For \(U'=\mathcal R_{a,A}U\), the associated isometry is
\begin{equation}\protect\phantomsection\label{eq:es-normal-isometry}{
 \Xi^U_{a,A}:g[U']\longrightarrow g[U],\qquad
 v\longmapsto
 \operatorname{Exp}^{g[U]}_{x_a}(\nu_aAv).
}\end{equation} Thus the instruction \(U\to U'\) corresponds to
\((\Xi^U_{a,A})^{-1}\). These maps, their inverses and the comparisons
in every conditioned finite diagram have common representatives for
bounded holomorphic record families, and depend holomorphically on the
retained parameters.

\emph{Proof.} Along \(\gamma_a(t)=\operatorname{Exp}^{g[U]}_o(ta)\),
read the free records \(Y_r(t)\) in the spin frame parallel-transported
from the mark. Every component is holomorphic in \(t\). Proposition
\hyperref[prop:es-completion]{3.2} identifies the ordered curvature
derivatives of this Einstein metric at each point with the universal
completion of its free record. Differentiation along \(\gamma_a\), whose
velocity has constant frame components \(a\), therefore gives
\begin{equation}\protect\phantomsection\label{eq:es-geometric-record-flow}{
 \frac{\mathrm d}{\mathrm dt}Y_r(t)
       =i_aY_{r+1}(t)+Q_{r,a}(Y(t)),\qquad Y_r(0)=U_r.
}\end{equation} The analytic-time coefficient recursion of Theorem
\hyperref[thm:es-flow]{6.1} determines every component from these
initial values. Thus \(Y(t)=\operatorname{Fl}_{\mathsf H_a}^{\,t}U\) on
the common time domain. Applying the frame change \(A\) at \(t=1\)
proves that the record at \(x_a\) in the frame \(\nu_aA\) is \(U'\).

Rebase \(g[U]\) at \(x_a\) and give it that frame. Proposition
\hyperref[prop:es-normal-uniqueness]{7.1} identifies this rebased metric
in its normal coordinates with \(g[U']\). The resulting isometry is
exactly eq.~\ref{eq:es-normal-isometry}: the new normal-coordinate
representation maps into the original one. Its differential at the
marked point is \(\nu_aA\), and its spin lift is the one obtained from
the initialized radial transport and the specified element \(A\).

Conversely, let \(f\) be such a local isometry from a neighborhood of
the mark of \(g[U']\) into a retained representative of \(g[U]\). Its
marked image determines \(a=(\operatorname{Exp}^{g[U]}_o)^{-1}(f(0))\).
Its specified spin differential, expressed relative to \(\nu_a\),
determines a unique full spin element \(A\). Curvature and covariant
differentiation are natural under an isometry, so their complete
symmetric records give \(U'=\mathcal R_{a,A}U\). An isometry carries
affinely parametrized geodesics to affinely parametrized geodesics and
preserves their initial tangent vectors. Hence it agrees with
eq.~\ref{eq:es-normal-isometry} on a normal neighborhood. This proves
both existence and uniqueness of the comparison.

Theorem \hyperref[thm:es-flow]{6.1} gives common record-flow domains;
Proposition \hyperref[prop:es-uniform-realization]{7.3} bounds the
metric, inverse metric and their finite derivatives on common normal
parents. The geodesic and parallel-spin equations are finite analytic
ODEs there. Their solutions and inverse exponential maps therefore have
common smaller domains and holomorphic parameter dependence. For a
conditioned finite diagram, the minimum of its finitely many positive
buffers and inverse margins gives common domains for all representatives
and specified composites. ∎

The record composition law compares a broken geodesic with the direct
radial geodesic to its endpoint. Define
\begin{equation}\protect\phantomsection\label{eq:es-endpoint-composition}{
 h_U(a,c)
 =\bigl(\operatorname{Exp}^{g[U]}_o\bigr)^{-1}
    \operatorname{Exp}^{g[U]}_{x_a}(\nu_ac).
}\end{equation} Let \(\nu^{a}_c\) be spin parallel transport along the
second geodesic, from \(x_a\) with initial tangent \(\nu_ac\), to the
endpoint in eq.~\ref{eq:es-endpoint-composition}. With \(h=h_U(a,c)\),
set \begin{equation}\protect\phantomsection\label{eq:es-spin-holonomy}{
 K_U(a,c)=\nu_h^{-1}\nu^{a}_c\nu_a.
}\end{equation} This is the full spin rotation comparing broken radial
transport with direct radial transport. Its value at a constant broken
path is the identity.

\pandocstatementheading{\protect\phantomsection\label{prop:es-composition}Proposition 8.2.}

For successive admissible comparisons \((a,A):U\to U'\) and
\((b,B):U'\to U''\), composition and inversion are
\begin{equation}\protect\phantomsection\label{eq:es-record-composition}{
 \begin{aligned}
 (b,B)\circ(a,A)
   &=\bigl(h_U(a,Ab),\ K_U(a,Ab)AB\bigr),\\
 (a,A)^{-1}
   &=\bigl(-A^{-1}a,A^{-1}\bigr)\quad\text{at }U'.
 \end{aligned}
}\end{equation} These operations define the local comparison groupoid
and identify it with the local spin-isometry groupoid under Theorem
\hyperref[thm:es-physical-arrows]{8.1}.

\emph{Proof.} The composite isometry \(\Xi^U_{a,A}\circ\Xi^{U'}_{b,B}\)
sends the final mark to \(\operatorname{Exp}^{g[U]}_{x_a}(\nu_aAb)\).
Its normal displacement from \(o\) is \(h_U(a,Ab)\). An isometry
intertwines the spin parallel transports initialized by its spin
differential. Consequently its differential at the final mark is the
broken transport \(\nu^{a}_{Ab}\nu_aAB\). Expressing this relative to
the direct radial transport \(\nu_{h_U(a,Ab)}\) gives the second
component of the first formula.

The reversed first radial geodesic, expressed in the frame \(\nu_aA\) at
\(x_a\), has initial tangent \(-A^{-1}a\). Transport along that reversed
segment cancels \(\nu_a\). The remaining frame comparison is \(A^{-1}\),
giving the inverse formula with its full spin lift. The identity is
\((0,1)\). Associativity follows from actual composition of the
isometries and the uniqueness in Theorem
\hyperref[thm:es-physical-arrows]{8.1}. The same uniqueness proves that
these formulas agree on overlapping representatives. ∎

The element \((-1,-1)\in G\) acts trivially on \(V\) and on every tensor
record. It remains a distinct specified lift in the comparison groupoid.
Thus the record comparison includes its spin label as well as its action
on tensor components.

Fix now the Lorentz real form \(V_{\mathbb R}\subset V\) of signature
\((+,-,-,-)\) and its spin conjugation. On records it induces the
antilinear involution \(\mathfrak c\) of Definition 2.4, which also
conjugates \(\Lambda\). Its fixed points in the fixed spin
identifications satisfy
\begin{equation}\protect\phantomsection\label{eq:es-lorentz-records}{
 \Lambda\in\mathbb R,\qquad
 U_r^-=\overline{U_r^+}\quad(r\ge0).
}\end{equation}

\pandocstatementheading{\protect\phantomsection\label{prop:es-lorentz-reality}Proposition 8.3.}

The fixed records in eq.~\ref{eq:es-lorentz-records} correspond exactly
to marked real-analytic Lorentzian Einstein four-germs. The
correspondence preserves the real point comparisons and full Lorentz
spin lifts of Theorem \hyperref[thm:es-physical-arrows]{8.1}.

\emph{Proof.} By Proposition 3.2, curvature completion commutes with
\(\mathfrak c\). Hence the horizontal and vertical derivations satisfy
\(\mathfrak c(\mathsf H_aX)=\mathsf H_{\mathfrak ca}(\mathfrak cX)\) and
the corresponding identity for spin-frame derivations. Uniqueness of the
convergent initialized flow gives
\begin{equation}\protect\phantomsection\label{eq:es-real-flow}{
 \mathfrak c\,\operatorname{Fl}_{\mathsf H_a}^{\,t}X
 =\operatorname{Fl}_{\mathsf H_{\mathfrak ca}}^{\,\bar t}
      (\mathfrak cX).
}\end{equation} Conjugate the normal metric by the fixed antilinear map
on \(V\): its value at \(x\) on \(v,w\) is
\(\overline{g(\mathfrak cx)(\mathfrak cv,\mathfrak cw)}\). This is a
holomorphic normal metric with the prescribed tangent form and conjugate
Einstein coefficient. Its free records are \(\mathfrak cX\). Proposition
\hyperref[prop:es-normal-uniqueness]{7.1} therefore identifies it with
\(g[\mathfrak cX]\). For fixed input, \(g[X]\) is consequently real on
\(V_{\mathbb R}\). At the mark it equals the prescribed nondegenerate
Lorentz form. Continuity and a restriction of the real neighborhood give
the same signature throughout, and its Einstein coefficient is the real
prescribed \(\Lambda\).

Conversely, a real-analytic Lorentzian Einstein metric has a local
holomorphic extension, obtained by complexifying its coordinate
coefficients and its marked spin frame. The Einstein equation extends by
the identity theorem, since it is an analytic identity in the metric,
its inverse and its first two derivatives. Its free curvature
derivatives satisfy eq.~\ref{eq:es-lorentz-records}. Their factorial
estimate follows from the converse part of Theorem
\hyperref[thm:es-realization]{7.2}. The normal realization with these
records is identified with the holomorphic extension by Proposition
\hyperref[prop:es-normal-uniqueness]{7.1}. Restriction to the real slice
gives the original real-analytic metric.

For real directions and real times, eq.~\ref{eq:es-real-flow} preserves
the fixed records. Real geodesic development and parallel transport give
the real point and spin maps. Conversely a real local isometry with a
prescribed Lorentz spin lift yields, by the extraction in Theorem
\hyperref[thm:es-physical-arrows]{8.1}, precisely a real displacement
and that real spin element. ∎

\protect\phantomsection\label{sec:es-finite-witnesses}

\subsection{Appendix A. Finite Einstein
witnesses}\label{appendix-a.-finite-einstein-witnesses}

The completion argument requires analytic realizations of finite record
prefixes. We construct them by solving analytic constraints and
evolution equations, then prescribing the prefix through a
finite-dimensional inverse theorem. We use the analytic Cauchy theorem
in the following normal form: a finite analytic system solved for its
highest derivatives transverse to a noncharacteristic analytic
hypersurface, with analytic Cauchy data, has a unique local analytic
solution. Additional undifferentiated finite-dimensional parameters give
a jointly holomorphic solution after restriction (Kowalevsky 1875).

\pandocstatementheading{\protect\phantomsection\label{lem:es-finite-witness}Lemma A.1.}

For every \(N\ge0\), every finite Lorentz-real prefix \(U_r^\pm\),
\(0\le r\le N\), and every real \(\Lambda\), there is a real-analytic
Lorentzian Einstein metric with this prefix, this numerical coefficient
and the prescribed tangent and spin marking.

\emph{Proof.} Choose a real orthonormal basis of the Lorentz tangent
space, so that the physical flat metric is
\(\eta=(\mathrm dt)^2-\sum_{i=1}^3(\mathrm dy^i)^2\). We first solve
analytic constraints and evolution with an auxiliary metric of signature
\((-,+,+,+)\) and coefficient \(\lambda\). The physical metric will be
its negative and will have coefficient \(\Lambda=-\lambda\). The inverse
theorem will first realize small prefixes; a final dilation will reach
arbitrary finite targets.

Let \(A_\nu,B_\nu\) be real-analytic symmetric spatial tensors
satisfying the flat transverse trace-free conditions
\begin{equation}\protect\phantomsection\label{eq:es-finite-tt-data}{
 \operatorname{tr}A_\nu=\operatorname{tr}B_\nu=0,
 \qquad \partial^j(A_\nu)_{ij}
       =\partial^j(B_\nu)_{ij}=0.
}\end{equation} Their finite-dimensional choice will be made below. For
real amplitude parameters \(c_\nu\) set
\(\bar\gamma=\delta+\sum_\nu c_\nu A_\nu\) and let \(\Sigma\) be the
\(\bar\gamma\)-trace-free part of \(\sum_\nu c_\nu B_\nu\). Put
\begin{equation}\protect\phantomsection\label{eq:es-finite-conformal-killing}{
 (\bar LZ)_{ij}
 =\bar\nabla_iZ_j+\bar\nabla_jZ_i
    -\frac23\bar\gamma_{ij}\bar\nabla_kZ^k .
}\end{equation} Solve the two analytic systems
\begin{equation}\protect\phantomsection\label{eq:es-finite-constraints}{
 \overline{\operatorname{div}}\bar LZ
       =-\overline{\operatorname{div}}\Sigma,
 \qquad
 \bar\Delta\phi-\frac18R(\bar\gamma)\phi
  +\frac18|\Sigma+\bar LZ|_{\bar\gamma}^2\phi^{-7}
  +\frac{\lambda}{4}\phi^5=0
}\end{equation} with \(Z=\partial_1Z=0\), \(\phi=1\),
\(\partial_1\phi=0\) on \(y^1=0\). Here tensor contractions are bilinear
in the complexified systems. At \(c=\lambda=0\), the coefficient of
\(\partial_1^2 Z\) in the first system is
\(I+\frac13(\mathrm dy^1)^\sharp\otimes\mathrm dy^1\), with eigenvalues
\(1,1,4/3\). The coefficient of \(\partial_1^2\phi\) in the second
system is \(\bar\gamma^{11}=1\). The first system is independent of
\(\phi\), and the second is analytic near the unit \(\phi=1\). The
analytic Cauchy theorem therefore gives solutions jointly holomorphic in
\(c,\lambda\) and \(y\) near zero.

Set \(T=\Sigma+\bar LZ\), \(\gamma=\phi^4\bar\gamma\) and
\(K=\phi^{-2}T\). The conformal connection difference is
\(2\phi^{-1}(\delta_i^k\partial_j\phi+\delta_j^k\partial_i\phi-\bar\gamma_{ij}\bar\gamma^{k\ell}\partial_\ell\phi)\).
Substitution gives
\begin{equation}\protect\phantomsection\label{eq:es-finite-conformal-identities}{
 \begin{aligned}
 R(\gamma)&=\phi^{-5}\bigl(-8\bar\Delta\phi+R(\bar\gamma)\phi\bigr),\\
 \operatorname{div}_\gamma K
   &=\phi^{-6}\overline{\operatorname{div}}T,&
 |K|_\gamma^2&=\phi^{-12}|T|_{\bar\gamma}^2,&
 \operatorname{tr}_\gamma K&=0.
 \end{aligned}
}\end{equation} Equations (\ref{eq:es-finite-constraints}) imply the
constraint equations
\begin{equation}\protect\phantomsection\label{eq:es-finite-solved-constraints}{
 R(\gamma)+(\operatorname{tr}_\gamma K)^2-|K|_\gamma^2
       =2\lambda,
 \qquad
 \operatorname{div}_\gamma K
       -\mathrm d\operatorname{tr}_\gamma K=0.
}\end{equation} At the flat parameter, the first amplitude variations of
\(Z\) and \(\phi-1\) vanish. Indeed their sources are respectively
\(\partial^j(B_\nu)_{ij}\) and
\(\partial_i\partial_j(A_\nu)_{ij}-\Delta\operatorname{tr}A_\nu\), and
their Cauchy data are zero. Thus \(\dot\gamma=A_\nu\) and
\(\dot K=B_\nu\) in each amplitude direction.

Write \(\tau=\operatorname{tr}_\gamma K\) and evolve by
\begin{equation}\protect\phantomsection\label{eq:es-finite-gaussian-evolution}{
 \partial_t\gamma_{ij}=2K_{ij},\qquad
 \partial_tK_{ij}
   =-\operatorname{Ric}(\gamma)_{ij}
     -\tau K_{ij}+2K_i{}^kK_{kj}
     +\lambda\gamma_{ij}.
}\end{equation} After eliminating \(K\), these equations are solved for
\(\partial_t^2\gamma\), with right-hand sides analytic in the invertible
spatial metric, its tangential derivatives through order two and its
first time derivative. The analytic initial data are the \(\gamma,K\)
constructed above, so the analytic Cauchy theorem applies with normal
coefficient the identity. Put
\(g_{\mathrm a}=-(\mathrm dt)^2+\gamma_{ij}\,\mathrm dy^i\mathrm dy^j\).

The time-containing Christoffel symbols are \(\Gamma^0_{ij}=K_{ij}\) and
\(\Gamma^i_{0j}=K^i{}_j\). Consequently
\begin{equation}\protect\phantomsection\label{eq:es-finite-gaussian-ricci}{\begin{aligned}
\operatorname{Ric}(g_{\mathrm a})_{ij}
 &=\operatorname{Ric}(\gamma)_{ij}
     +\partial_tK_{ij}+\tau K_{ij}-2K_i{}^kK_{kj}, \\
\operatorname{Ric}(g_{\mathrm a})_{00}&=-\partial_t\tau-|K|^2, \\
\operatorname{Ric}(g_{\mathrm a})_{0i}&=\nabla^jK_{ij}-\partial_i\tau.
\end{aligned}}\end{equation} The spatial Einstein equations hold by construction. Let
\(\mathcal H=\operatorname{Ric}(g_{\mathrm a})_{00}+\lambda\) and
\(\mathcal M_i=\operatorname{Ric}(g_{\mathrm a})_{0i}\). The contracted
Bianchi identity, with those spatial equations imposed, is the
homogeneous normal system
\begin{equation}\protect\phantomsection\label{eq:es-finite-constraint-propagation}{
 \partial_t\mathcal H
 =2\nabla^i\mathcal M_i-2\tau\mathcal H,
 \qquad
 \partial_t\mathcal M_i
 =\frac12\partial_i\mathcal H-\tau\mathcal M_i.
}\end{equation} The initial values vanish by
eq.~\ref{eq:es-finite-solved-constraints}. Analytic uniqueness gives
\(\mathcal H=\mathcal M=0\) and hence
\(\operatorname{Ric}(g_{\mathrm a})=\lambda g_{\mathrm a}\). The
initialized solutions are real on the real parameter and coordinate
slice by conjugation and uniqueness. On a smaller real neighborhood
\(\gamma\) is positive definite. Therefore \(g=-g_{\mathrm a}\) is
Lorentzian, and \(\operatorname{Ric}(g)=-\lambda g\).

It remains to prescribe the free record prefix. We choose a finite TT
family whose physical free-record map is a submersion with equal source
and target dimensions. For \(z\in\mathbb C\) put \(\kappa(z)=(1,z)\) and
\(k(z)_{AA'}=\kappa_A(z)\bar\kappa_{A'}(z)\). These are nonzero real
null covectors. For each \(r\), the tensors
\begin{equation}\protect\phantomsection\label{eq:es-finite-coherent-directions}{
 \kappa(z)^{r+4}\bar\kappa(z)^r
}\end{equation} span \(W_r^+\) over \(\mathbb C\). To see this, a
functional annihilating all of them gives a polynomial in \(z,\bar z\)
of bidegree at most \((r+4,r)\). With \(z=x+\mathrm i y\), it becomes a
polynomial vanishing for every real pair \((x,y)\), hence is identically
zero. The substitution between \((z,\bar z)\) and \((x,y)\) is
invertible over \(\mathbb C\), so the original functional vanishes. A
finite spanning set valid for every \(r\le N\) is obtained from
\(x,y\in\{0,\ldots,2N+4\}\): its polynomial degree in each real variable
is at most \(2N+4\), and successive one-variable interpolation proves
the assertion on this grid. Choose a complex basis of \((r+5)(r+1)\)
such directions for each \(r\).

Choose distinct nonzero real numbers \(\tau_0,\ldots,\tau_N\). The
Vandermonde matrix \(((\mathrm i\tau_\ell)^j)_{0\le j,\ell\le N}\) is
invertible. Its inverse columns give profiles
\begin{equation}\protect\phantomsection\label{eq:es-finite-profiles}{
 f_r(s)=\sum_{\ell=0}^N b_\ell^{(r)}e^{\mathrm i\tau_\ell s},
 \qquad f_r^{(j)}(0)=\delta_{jr}\quad(0\le j\le N).
}\end{equation} For one chosen direction write
\(k_{AA'}=\kappa_A\lambda_{A'}\) with \(\lambda=\bar\kappa\), and choose
\(\mu\) with \(d=\lambda_{0'}\mu_{1'}-\lambda_{1'}\mu_{0'}\ne0\). The
polarization \(p_{AA'}=\kappa_A\mu_{A'}\) is null and transverse to
\(k\). In the physical flat tangent convention the metric variation
\(h_{ab}=a p_ap_b e^{\mathrm i\tau k\cdot x}\) has linearized last-slot
curvature
\begin{equation}\protect\phantomsection\label{eq:es-finite-wave-readout}{
 \begin{aligned}
 R^{(1)}_{abcd}(h)
 &=\frac12\bigl(
   \partial_a\partial_ch_{bd}+\partial_b\partial_dh_{ac}
   -\partial_a\partial_dh_{bc}-\partial_b\partial_ch_{ad}\bigr)\\
 &=-\frac12a\tau^2
   (k_ap_b-k_bp_a)(k_cp_d-k_dp_c)e^{\mathrm i\tau k\cdot x},\\
 \Psi^{+(1)}_{ABCD}
 &=\frac12a\tau^2d^2\kappa_A\kappa_B\kappa_C\kappa_D
       e^{\mathrm i\tau k\cdot x},
 &\Psi^{-(1)}&=0.
 \end{aligned}
}\end{equation} The first line follows by differentiating the
Christoffel formula at a constant metric. For the last line use
\(k_ap_b-k_bp_a=d\kappa_A\kappa_B\epsilon_{A'B'}\) and the curvature
readout of eq.~\ref{eq:es-curvature-decomposition}. All normalization
factors displayed here are nonzero.

Replace \(p\) by \(p-(p_0/k_0)k\). Since \(k\) is a nonzero real null
covector, \(k_0\ne0\). This change removes \(p_0\) and preserves its
wedge with \(k\), its nullness and its transversality. The resulting
wave is temporal and its spatial initial metric and half time derivative
satisfy eq.~\ref{eq:es-finite-tt-data}. Pair it with its Lorentz
conjugate to obtain a real TT wave. Use the negative of this real wave
for the auxiliary initial data, so that the physical metric variation
after \(g=-g_{\mathrm a}\) is the displayed wave and its conjugate.

For each frequency divide its amplitude by \(\tau_\ell^2d^2/2\) and
combine with eq.~\ref{eq:es-finite-profiles}. The positive linearized
Weyl spinor is then \(c\kappa^4f_r(k\cdot x)\), its negative spinor is
its Lorentz conjugate, and its order-\(j\) free record at the origin is
\begin{equation}\protect\phantomsection\label{eq:es-finite-wave-record}{
 \delta_{jr}c\kappa^{r+4}\bar\kappa^r
 \qquad(0\le j\le N).
}\end{equation} Real and imaginary parts of \(c\) supply two real
amplitude directions. Put \(d_N=2\sum_{r=0}^N(r+5)(r+1)\). The chosen
bases and profiles therefore give \(d_N\) real amplitudes and an
invertible derivative onto the real space of Lorentz-real prefixes.

Read the nonlinear family in an analytic orthonormal spin frame
\(\nu(c,\lambda)\) at the origin, identified with the prescribed
marking. Such a frame and its local spin lift are obtained by analytic
orthonormalization near the flat frame. At the flat parameter, all
background Weyl derivatives vanish, so point-frame and connection
variations contribute nothing to the first variation of the free-record
map. Linearization of eq.~\ref{eq:es-finite-gaussian-evolution} has
precisely the TT wave with its prescribed initial data, by analytic
uniqueness. The analytic map to the real prefix space is
\begin{equation}\protect\phantomsection\label{eq:es-prefix-submersion}{
 \begin{aligned}
 \mathcal F_N:\mathbb R^{d_N}\times\mathbb R
    &\longrightarrow
    \left(\bigoplus_{r=0}^N(W_r^+\oplus W_r^-)^{\mathfrak c}\right)
       \times\mathbb R,\\
 (c,\lambda)&\longmapsto
    \bigl((U_r(g(c,\lambda),\nu(c,\lambda)))_{r=0}^N,-\lambda\bigr),\\
 D\mathcal F_N(0,0)&=
       \begin{pmatrix}L_N&*\\0&-1\end{pmatrix},
       \qquad L_N\text{ a real linear isomorphism}.
 \end{aligned}
}\end{equation} The finite-dimensional analytic inverse theorem realizes
every sufficiently small real prefix and scalar. For an arbitrary finite
target choose \(s>0\) so large that
\((s^{-r-2}U_r,s^{-2}\Lambda)_{0\le r\le N}\) lies in this neighborhood.
If \(h\) realizes the scaled target, set \(g=s^{-2}\delta_s^*h\) with
\(\delta_s(x)=sx\). The pullback contributes \(s^{r+4}\) on the \(r+4\)
covariant curvature-derivative slots, and the constant metric scaling
contributes \(s^{-2}\). Thus the tangent metric is unchanged, the
order-\(r\) marked curvature derivative is multiplied by \(s^{r+2}\),
and the Einstein coefficient by \(s^2\). It realizes the requested
prefix and coefficient. ∎

This construction uses finitely many analytic equations and
finite-dimensional parameters. Its neighborhood may depend on the length
and values of the prescribed prefix. Proposition
\hyperref[prop:es-completion]{3.2} uses this finite-prefix existence to
verify polynomial identities; common analytic domains are obtained
separately from the record-flow and Jacobi estimates.

\protect\phantomsection\label{sec:es-low-orders}

\subsection{Appendix B. The first nonlinear
corrections}\label{appendix-b.-the-first-nonlinear-corrections}

The first three corrections show how linear shifts lead to the first
nonlinear insertion. Recursion (\ref{eq:es-recursion}) gives finite
tensor formulas at each fixed order, using the same compatible Spencer
inverse \(M_rJ_r^*\) as in §4. For a polynomial coefficient \(F_b\) let
\(D_jF_b[Z]\) mean its differential with respect to the two chiral
components of \(U_j\); the scalar \(\Lambda\) is held fixed. The two
outer derivative arguments in the following formulas are \(a,b\).

Since \(Q_0=0\), the first correction is
\begin{equation}\protect\phantomsection\label{eq:es-q-one}{
 Q_1=-M_1J_1^*K_1,\qquad
 K_1(a,b)=\mathsf V_{R_U(a,b)}U_0.
}\end{equation} It is quadratic in \((U_0,\Lambda)\). Define the
two-argument tensor
\begin{equation}\protect\phantomsection\label{eq:es-q-two}{\begin{gathered}
(T_2)_{a,b}=D_0Q_{1,b}[i_aU_1], \\
Q_2=M_2J_2^*(\tau T_2-K_2), \\
K_2(a,b)=\mathsf V_{R_U(a,b)}U_1.
\end{gathered}}\end{equation} Every term of \(Q_2\) has curvature degree two and total
input order one. At the next order put
\begin{equation}\protect\phantomsection\label{eq:es-q-three}{
 \begin{aligned}
 (T_3^{[2]})_{a,b}
     &=D_0Q_{2,b}[i_aU_1]+D_1Q_{2,b}[i_aU_2],\\
 (T_3^{[3]})_{a,b}
     &=D_1Q_{2,b}[Q_{1,a}],\\
 Q_3^{[2]}&=M_3J_3^*(\tau T_3^{[2]}-K_3),\\
 Q_3^{[3]}&=M_3J_3^*\tau T_3^{[3]},\qquad
 Q_3=Q_3^{[2]}+Q_3^{[3]}.
 \end{aligned}
}\end{equation} Here \(K_3(a,b)=\mathsf V_{R_U(a,b)}U_2\). The quadratic
inputs have total order two, while the cubic inputs all have order zero.
The last line exhibits the first possible nonlinear insertion: a
coefficient of \(Q_1\) replaces the order-one position of \(Q_2\). Both
output chiralities and both possible intermediate record chiralities are
included in the differential \(D_1\). Differentiation of a scalar
position contributes zero.

These formulas follow by applying the polynomial product rule in
eq.~\ref{eq:es-recursion}. The weight relation
eq.~\ref{eq:es-completion-weight} permits only degree two for \(r=1,2\)
and degrees two and three for \(r=3\), so the displayed formulas contain
every term at these orders. The bounds for all orders are given by
Theorem \hyperref[thm:es-coefficients]{5.5}.

The recursion can also be checked directly on the Ricci-flat Kasner
metric \[
 g=\mathrm dt^2-\sum_{i=1}^3t^{2p_i}(\mathrm dx^i)^2,
 \qquad (p_1,p_2,p_3)=(-2/7,3/7,6/7).
\] In its orthonormal frame, compute the ordered curvature derivatives
through order four at \(t=1\) and project to the two chiral records.
Evaluate \(Q_r=\nabla U_r-iU_{r+1}\) and \(\mathsf H Q_{r-1}\) by
covariant differentiation of these geometric records, and \(K_r\) by
curvature action. Applying the Spencer inverse to the right side of
eq.~\ref{eq:es-recursion} agrees with \(Q_r\) in all 520 components for
\(1\le r\le3\), with zero fully symmetric projection. These are exact
equalities over \(\mathbb Q(\mathrm i)\) after a common spin-basis
rescaling. The companion archive,
\href{exact-set-checks-v0.1.1.zip}{\texttt{exact-set-checks-v0.1.1.zip}},
is included with the source and computational supplements and supplies
these calculations, their exact output and checksums. Its SageMath 10.9
entry point is \texttt{check\_kasner\_recursion.py\ -\/-max-r\ 3}; the
accompanying README specifies the spin-basis rescaling, curvature
convention and isolated execution environment.

\protect\phantomsection\label{sec:es-ai-use}

\subsection*{AI-use disclosure}\label{ai-use-disclosure}

AI systems were used extensively throughout this research and the
preparation of this manuscript. The models used were GPT-6 Pro,
gpt-6.1-sol (xhigh), and Claude Opus 5.5 (xhigh).

\section*{References}

\protect\phantomsection\label{refs}
\begin{CSLReferences}{1}{1}
\sloppy
\bibitem[\citeproctext]{ref-Acena2009}
Aceña, Andrés E. 2009. {``Convergent Null Data Expansions at Space-Like
Infinity of Stationary Vacuum Solutions.''} \emph{Annales Henri
Poincar{é}} 10: 275--337.
\url{https://doi.org/10.1007/s00023-009-0406-z}.

\bibitem[\citeproctext]{ref-Acena2010}
Aceña, Andrés E. 2010. {``Minimal Data at a Given Point of Space for
Solutions to Certain Geometric Systems.''} \emph{Classical and Quantum
Gravity} 27 (15): 155006.
\url{https://doi.org/10.1088/0264-9381/27/15/155006}.

\bibitem[\citeproctext]{ref-AndersonTorre1996}
Anderson, Ian M., and Charles G. Torre. 1996. {``Classification of Local
Generalized Symmetries for the Vacuum {Einstein} Equations.''}
\emph{Communications in Mathematical Physics} 176 (3): 479--539.
\url{https://doi.org/10.1007/BF02099248}.

\bibitem[\citeproctext]{ref-BasurtoArzate2026}
Basurto-Arzate, Efrain. 2026. {``Torsion-Free Connections with
Prescribed Curvature.''} \emph{Advances in Geometry} 26 (3): 387--406.
\url{https://doi.org/10.1515/advgeom-2026-0011}.

\bibitem[\citeproctext]{ref-ChruscielPaetz2013}
Chruściel, Piotr T., and Tim-Torben Paetz. 2013. {``Solutions of the
Vacuum {Einstein} Equations with Initial Data on Past Null Infinity.''}
\emph{Classical and Quantum Gravity} 30 (23): 235037.
\url{https://doi.org/10.1088/0264-9381/30/23/235037}.

\bibitem[\citeproctext]{ref-Friedrich2007}
Friedrich, Helmut. 2007. {``Static Vacuum Solutions from Convergent Null
Data Expansions at Space-Like Infinity.''} \emph{Annales Henri
Poincar{é}} 8: 817--84. \url{https://doi.org/10.1007/s00023-006-0323-3}.

\bibitem[\citeproctext]{ref-Friedrich2013}
Friedrich, Helmut. 2013. {``The {Taylor} Expansion at Past Time-Like
Infinity.''} \emph{Communications in Mathematical Physics} 324 (1):
263--300. \url{https://doi.org/10.1007/s00220-013-1803-1}.

\bibitem[\citeproctext]{ref-Gasqui1982}
Gasqui, Jacques. 1982. {``Sur la r{é}solubilit{é} locale des {é}quations
d'{Einstein}.''} \emph{Compositio Mathematica} 47 (1): 43--69.
\url{https://www.numdam.org/item/CM_1982__47_1_43_0/}.

\bibitem[\citeproctext]{ref-Isthmus2026Twistor}
Isthmus, Lior. 2026. \emph{A Projective Twistor Correspondence for Local
Analytic {Einstein} Germs}. Hexagon preprint hexagon:2611.00123v1.
\url{https://hexagonmath.org/2611.00123v1}.

\bibitem[\citeproctext]{ref-Jentsch2026}
Jentsch, Tillmann. 2026. \emph{The Jet Isomorphism Theorem of
{Riemannian} Geometry}. \url{https://arxiv.org/abs/1509.08269v2}.

\bibitem[\citeproctext]{ref-Kowalevsky1875}
Kowalevsky, Sophie von. 1875. {``{Zur Theorie der partiellen
Differentialgleichungen}.''} \emph{{Journal für die reine und angewandte
Mathematik}} 80: 1--32.

\bibitem[\citeproctext]{ref-KowalskiBelger1994}
Kowalski, Oldřich, and Martin Belger. 1994. {``{Riemannian} Metrics with
the Prescribed Curvature Tensor and All Its Covariant Derivatives at One
Point.''} \emph{Mathematische Nachrichten} 168 (1): 209--25.
\url{https://doi.org/10.1002/mana.19941680113}.

\bibitem[\citeproctext]{ref-Penrose1980}
Penrose, Roger. 1980. {``Golden Oldie: Null Hypersurface Initial Data
for Classical Fields of Arbitrary Spin and for General Relativity.''}
\emph{General Relativity and Gravitation} 12 (3): 225--64.
\url{https://doi.org/10.1007/BF00756234}.

\bibitem[\citeproctext]{ref-Torre1995}
Torre, Charles G. 1995. \emph{Spinors, Jets, and the {Einstein}
Equations}. \url{https://arxiv.org/abs/gr-qc/9508005}.

\end{CSLReferences}

\end{document}
